% concatenated from GuZheng_Final_Version1_20251213.tex
\documentclass{article}

\usepackage{amsmath}
\usepackage{mathrsfs}
\usepackage{amssymb}
\usepackage{amsthm}

\usepackage{authblk}

\usepackage{geometry}
\geometry{scale=0.7}

\newtheorem{lemma}{Lemma}[section]
\newtheorem{theorem}[lemma]{Theorem}
\newtheorem{corollary}[lemma]{Corollary}
\newtheorem{proposition}[lemma]{Proposition}
\newtheorem{definition}[lemma]{Definition}
\newtheorem{criterion}[lemma]{Criterion}
\newtheorem{example}[lemma]{Example}
\newtheorem{remark}[lemma]{Remark}
\newtheorem{problem}[lemma]{Problem}

\newcommand{\cc}{{\mathcal C}}
\def\1{{\bf 1}}
\newcommand{\id}{{\rm id}}
\newcommand{\Hom}{{\rm Hom}\,}
\newcommand{\ot}{\otimes}

\newcommand{\ito}[3]{\left(^{\ \ {#3}}_{{#1}\ {#2}}\right)}

\begin{document}
	\title{Uniqueness of vertex operator algebras
		arising from GKO-construction}
	
	\author{Yuhan Gu}
	\affil{School of Mathematics and Statistics, 
		Qingdao University}
	\author{Wen Zheng
		\footnote{Corresponding author\\
			Email address: $wzheng14@qdu.edu.cn$\\
			The author is supported by the NSFC No.~12201334 and the National Science Foundation of Shandong Province No.~ZR2022QA023.}}
	\affil{School of Mathematics and Statistics,
		Qingdao University}
	
	\maketitle
	
	\begin{abstract}
		A series of vertex operator algebras are constructed 
		by 
		% picking modules of kown vertex operator algebras and making them 
		% a mirror extension with certain summands in 
		GKO-construction, which is a generalization of 3A-algebra and 6A-algebra. 
		It is proved their vertex operator algebra structures are unique under nonzero assumptions on some elements of braiding matrices.
		Furthermore, we show each of them is generated by weight two subspace, i.e. the Griess algebra.
		% Braiding submatrices appear in the construction
		% can be computed as a side effect.
	\end{abstract}
	
	\setcounter{tocdepth}{2}
	\tableofcontents
	
	
	
	\newpage
	
	\section{Introduction}
	Constructing a series of vertex operator algebras (VOAs) $\mathcal{U}_k$ by GKO construction,
	and proving that $\mathcal{U}_k$ has a unique VOA structure
	are the main missions of this paper. 
	We briefly chart the history about structures and classifications of VOA
	before explaining any detailed methodology.
	
	Vertex (operator) algebra was introduced and studied
	extensively in mathematics 
	since 1980s \cite{B,FLM,FHL}. 
	Its properties show analogies of 
	commutative algebras, associative algebras,
	differential algebras and graded algebras.
	With these enriched structures,
	vertex operator algebras have built 
	deep connection between various aspects such as 
	group theories, modular forms, Lie algebras, 
	conformal nets, 
	tensor categories, and mathematical physics, etc,
	making it one of the most colorful field of study.
	
	People have already known a lot of examples of VOAs,
	such as Virasoro VOAs, Heisenberg VOAs, Affine VOAs,
	and Lattice VOAs. Apart from exsiting examples,
	new VOAs can be constructed by extensions, coset constructions,
	and orbifold theories.
	
	One of the most important VOAs 
	is the moonshine VOA $V^{\natural}$
	arising from the study of moonshine conjecture.
	It is proved in \cite{DMZ} that $V^{\natural}$ contains a 
	subalgebra $L(1/2,0)^{\otimes 48}$,
	where $L(1/2,0)$ is the simple Virasoro VOA with central charge $1/2$.
	
	Virasoro VOAs, especially unitary series of Virasoro VOAs, are the simplest and well-studied VOAs, which
	have many ideal properties that clarify the structure of them.
	Especially, we use GKO-construction 
	to get a series of VOAs $\mathcal{U}_k$, 
	and the fusion rules and braiding matrices \cite{FFK,H1,H2} of unitary series of Virasoro VOAs
	paly a vital role in our computations.
	
	One of our motivations is to generalize previous results about 3A-algebra \cite{M,SY} as well as 6A-algebra \cite{DJY,JZ}. For example, the $k=0$ and $k=1$ cases of VOAs $\mathcal{U}_k$ we constructed correspond  3A-algebra and 6A-algebra respectively. See Section 3 for the details. Another motivation is to try to understand Conjecture A.14 in \cite{LY}, which says the VOA generated by several Ising vectors related to 3A-algebra and 2A-algebra has unique VOA structure.
	In fact, VOA structures generated by two Ising vectors have been extensively studied \cite{M,SY,LYY, DJY, DZ, JZ, Zheng}. We would like to use similar methods like braiding matrices and fusion rules of Virasoro VOAs in \cite{DJY,DZ} to prove the uniqueness of VOA structure of $\mathcal{U}_k$, but the difficult part is to prove the uniqueness of $\mathcal{U}_k$ for all $k$ simultaneously.
	We give this proof under the nonvanishing assumption of some elements of the braiding matrices. For explicit computations of the elements of braiding matrices of $\mathcal{U}_1$ and $\mathcal{U}_2$, See Example 5.9 and 5.10.
	 
	%Bunches of new methods have been developed,
	%integrated with many other approaches, 
	%and yielded fruitful outputs such as 
	%tensor category theories \cite{H1,H2,H3},
	%quantum dimension theories \cite{DJX},
	%mirror extension theories \cite{X1,DJX2,Lin},
	%gluing methods \cite{CKM}, and
	%minimal modular extensions \cite{DNR}.
	%We will prove several fundamental results 
	%by directly using mirror extension theories.
	
	%The study of the uniqueness of VOA structure started years ago by \cite{DJY,DZ,JZ}.
	%Since they are deeply connected to subalgebras of the moonshine VOA,
	%new discoveries are found by using these propositions.
	%It is used by
	%\cite{Z2} to exploring the structure of VOAs generated by two Ising vectors.
	%It is also used to find the automorphism generated
	%by $\sigma$-involutions in parafermion vertex operator algebras
	%in \cite{LY}.
	%Some examples are clarified by uniqueness when
	%classifying bosonic rational conformal theories
	%at genera with $\operatorname{rank}(\mathcal{C})>4$ by \cite{R}.
	
	Now we give the sketch of our work. 
	By directly applying GKO-construction,
	we can decompose 
	$$\left(\mathcal{L}\left(3,0\right)\otimes\mathcal{L}\left(1,0\right)\otimes\mathcal{L}\left(1,0\right)\otimes\mathcal{L}\oplus\mathcal{L}\left(3,3\right)\otimes\mathcal{L}\left(1,1\right)\otimes\mathcal{L}\left(1,0\right)\right)\otimes\mathcal{L}(1,0)^{\otimes k}$$
	(See Section 2.5 for details) into
	$\bigoplus_{i}\mathcal{U}_{k,i}\otimes \mathcal{L}(k+5,i)$.
	$\mathcal{U}:=\mathcal{U}_k:=\mathcal{U}_{k,0}$ is defined by taking the commutant of $\mathcal{L}(k+5,0)$ in GKO-construction.%
	%, and it is a VOA.
	% where $\mathcal{U}_{k,i}$ are known VOA-modules and $\mathcal{U}:=\mathcal{U}_k:=\mathcal{U}_{k,0}$.
	%Let $\mathcal{U}_0$ be the $3A$-algebra, and it appears in the GKO-construction.
	More precisely, we have the decomposition
	$\mathcal{U}_k=\bigoplus U^i=\bigoplus P^i\otimes Q^i$,
	where $P_i$ are $\mathcal{U}_{k-1}$ modules and $Q_i$ are modules of Virasoro VOA. 
	%It is known in \cite{DJY} that $\mathcal{U}_{1}$ is the $6A$-algebra.
	% Given a VOA $\mathcal{U}_{k-1}$ which satisfies necessary conditions,
	% we pair its certain modules $P_i$ with Virasoro VOA models $Q_i$ appear
	% as summand before $\mathcal{L}(c_{k+6},0)$ in the 
	% GKO-construction, and make a mirror extension
	% by direct sums of tensor products.
	% That is, let $U^i:=P^i\otimes Q^i$,
	% we define $\mathcal{U}_k:=\bigoplus U^i$.
	Thus possible vertex operator is determined by fusion rules between those modules.
	We proved that fusion rules for $Q_i$ are enough for us.
	Since fusion rules for $Q_i$ are either $1$ or $0$, 
	we can put $Y=\sum \lambda_{i,j}^k \mathcal{Y}_{i,j}^k$,
	where $\mathcal{Y}_{i,j}^k$ is basis of intertwining operator of type 
	$\ito{U^i}{U^j}{U^k}$.
	%Since $\lambda_{i,j}^k=1$ for all $i,j,k$ compatible with fusion rules make a VOA structure $(\mathcal{U}, Y)$,
	Assume we fix $Y$ such that $\lambda_{i,j}^k=1$, i.e. $Y=\sum \mathcal{Y}_{i,j}^k$. For another possible VOA structure $\overline{Y}=\sum \lambda_{i,j}^k \mathcal{Y}_{i,j}^k$,
	we need to show $(\mathcal{U},Y)\cong(\mathcal{U},\overline{Y})$.
	Using mirror extension theory \cite{Lin}, we prove $\lambda_{i,j}^k\neq 0$ for all $i,j,k$ compatible with fusion rules.
	Under the nonvanishing assumption of some elements of the braiding matrices,
	we can assert $\left(\lambda_{i,j}^k\right)^2=1$ if $N_{i,j}^k=1$.
	Then we use $\lambda_{i,j}^k$ to define a map $\sigma$ between $(\mathcal{U},Y)$ and $(\mathcal{U},\overline{Y})$,
	and prove it is exactly the isomorphism we want.
	Furthermore, we show that $\mathcal{U}$ is generated by Griess algebra.
	
	This paper is organized as follows: 
	In Section 2, we recall some basic facts about
	vertex operator algebras and related topics needed in this paper. 
	In Section 3, the module $\mathcal{U}_k$ is constructed.
	In Section 4, we prove $\lambda_{i,j}^k\neq 0$ for all $i,j,k$ compatible with fusion rules.
	In Section 5, the isomorphism $\sigma$ is defined, and our main goal, the uniqueness of VOA structure on $\mathcal{U}_k$, is proved.
	In Section 6, we prove that $\mathcal{U}_k$ are generated by Griess algebra.
	In Section 7, we give some open problems based on this paper and its citations.
	
	\section{Preliminary}
	In this section, we review the basics on vertex operator algebras,
	the theory of quantum dimensions from \cite{DJX}, the coset realization
	of the discrete series of the unitary representations for the Virasoro
	algebra \cite{GKO}, the braiding matrices for certain Virasoro
	vertex operator algebras \cite{FFK}, 
	the structure of $6A$-algebra \cite{DJY},
	and mirror extensions for rational vertex operator algebras \cite{Lin}.
	
	\subsection{Basics}
	Let $V=(V,\ Y,\ \mathbf{1},\ \omega)$ be a vertex operator algebra.
	Let $Y(v,\ z)=\sum_{n\in\mathbb{Z}}v_{n}z^{-n-1}$ denote the vertex
	operator of $V$ for $v\in V$, where $v_{n}\in\mbox{End}(V)$. We
	first recall some basic notions from \cite{FLM,Z,DLM1,DLM3}.
	
	\begin{definition} $V_{2}$ is called \emph{Griess algebra}. A vector $v\in V_{2}$ is called a \emph{Virasoro
			vector with central charge $c_{v}$ }if it satisfies\emph{ $v_{1}v=2v$
		}and $v_{3}v=\frac{c_{v}}{2}\mathbf{1}$. Then the operators $L_{n}^{v}:=v_{n+1},\ n\in\mathbb{Z}$,
		satisfy the Virasoro commutation relation
		\[
		\left[L_{m}^{v},\ L_{n}^{v}\right]=\left(m-n\right)L_{m+n}^{v}+\delta_{m+n,\ 0}\frac{m^{3}-m}{12}c_{v}
		\]
		for $m,\ n\in\mathbb{Z}.$ A Virasoro vector $v\in V_{2}$ with central
		charge $1/2$ is called an \emph{Ising vector }if $v$ generates the
		Virasoro vertex operator algebra $L(1/2,\ 0)$.
		
	\end{definition}
	
	\begin{definition} An \emph{automorphism} $g$ of a vertex operator
		algebra $V$ is a linear isomorphism of $V$ satisfying $g(\bf{1})=\bf{1}$, $g\left(\omega\right)=\omega$
		and $gY\left(v,z\right)g^{-1}=Y\left(gv,z\right)$ for any $v\in V$.
		We denote by $\mbox{Aut}\left(V\right)$ the group of all automorphisms
		of $V$. \end{definition}
	
	For a subgroup $G\le\mbox{Aut}\left(V\right)$ the fixed point set
	$V^{G}=\left\{ v\in V|g\left(v\right)=v,\forall g\in G\right\} $
	has a vertex operator algebra structure.
	
	Let $g$ be an automorphism of a vertex operator algebra $V$ of order
	$T$. Denote the decomposition of $V$ into eigenspaces of $g$ as:
	
	\[
	V=\oplus_{r\in\mathbb{Z}/T\text{\ensuremath{\mathbb{Z}}}}V^{r}
	\]
	where $V^{r}=\left\{ v\in V|gv=e^{2\pi ir/T}v\right\} $.
	
	\begin{definition} A \emph{weak $g$-twisted $V$-module} $M$ is
		a vector space with a linear map
		\begin{align*}
			Y_{M}: & V\to\left(\text{End}M\right)\{z\}\\
			& v\mapsto Y_{M}\left(v,z\right)=\sum_{n\in\mathbb{Q}}v_{n}z^{-n-1}\ \left(v_{n}\in\mbox{End}M\right)
		\end{align*}
		which satisfies the following: for all $0\le r\le T-1$, $u\in V^{r}$,
		$v\in V$, $w\in M$,
		
		\[
		Y_{M}\left(u,z\right)=\sum_{n\in-\frac{r}{T}+\mathbb{Z}}u_{n}z^{-n-1},
		\]
		
		\[
		u_{l}w=0\ for\ l\gg0,
		\]
		
		\[
		Y_{M}\left(\mathbf{1},z\right)=Id_{M},
		\]
		
		\[
		z_{0}^{-1}\text{\ensuremath{\delta}}\left(\frac{z_{1}-z_{2}}{z_{0}}\right)Y_{M}\left(u,z_{1}\right)Y_{M}\left(v,z_{2}\right)-z_{0}^{-1}\delta\left(\frac{z_{2}-z_{1}}{-z_{0}}\right)Y_{M}\left(v,z_{2}\right)Y_{M}\left(u,z_{1}\right)
		\]
		
		\begin{equation}
			=z_{2}^{-1}\left(\frac{z_{1}-z_{0}}{z_{2}}\right)^{-r/T}\delta\left(\frac{z_{1}-z_{0}}{z_{2}}\right)Y_{M}\left(Y\left(u,z_{0}\right)v,z_{2}\right),\label{Jacobi for twisted V-module}
		\end{equation}
		where $\delta\left(z\right)=\sum_{n\in\mathbb{Z}}z^{n}$.
		
	\end{definition}
	
	\begin{definition}
		
		A $g$-\emph{twisted $V$-module} is a weak $g$-twisted $V$-module
		$M$ which carries a $\mathbb{C}$-grading induced by the spectrum
		of $L(0)$ where $L(0)$ is the component operator of $Y(\omega,z)=\sum_{n\in\mathbb{Z}}L(n)z^{-n-2}.$
		That is, we have $M=\bigoplus_{\lambda\in\mathbb{C}}M_{\lambda},$
		where $M_{\lambda}=\left\{ w\in M|L(0)w=\lambda w\right\} $. Moreover,
		$\dim M_{\lambda}$ is finite and for fixed $\lambda,$ $M_{\frac{n}{T}+\lambda}=0$
		for all small enough integers $n.$
		
	\end{definition}
	
	\begin{definition}An \emph{admissible $g$-twisted $V$-module} $M=\oplus_{n\in\frac{1}{T}\mathbb{Z}_{+}}M\left(n\right)$
		is a $\frac{1}{T}\mathbb{Z}_{+}$-graded weak $g$-twisted module
		such that $u_{m}M\left(n\right)\subset M\left(\mbox{wt}u-m-1+n\right)$
		for homogeneous $u\in V$ and $m,n\in\frac{1}{T}\mathbb{Z}.$
		
	\end{definition}
	
	If $g=Id_{V}$ we have the notions of weak, ordinary and admissible
	$V$-modules \cite{DLM2}.
	
	\begin{definition}A vertex operator algebra $V$ is called \emph{$g$-rational}
		if the admissible $g$-twisted module category is semisimple. $V$
		is called \emph{rational} if $V$ is $1$-rational. \end{definition}
	
	It was proved in \cite{DLM2} that if $V$ is rational then there
	are only finitely irreducible admissible $V$-modules up to isomorphism
	and each irreducible admissible $V$-module is ordinary. Let $M^{0},M^{1},$
	$\cdots,M^{d}$ be all the irreducible modules up to isomorphism with
	$M^{0}=V$. Then there exists $h_{i}\in\mathbb{C}$ for $i=0,\cdots,d$
	such that
	\[
	M^{i}=\oplus_{n=0}^{\infty}M_{h_{i}+n}^{i}
	\]
	where $M_{h_{i}}^{i}\not=0$ and $L\left(0\right)|_{M_{h_{i}+n}^{i}}=h_{i}+n$,
	$\forall n\in\mathbb{Z}_{+}$. $h_{i}$ is called the \emph{conformal
		weight} of $M^{i}$. We denote $M^{i}\left(n\right)=M_{h_{i}+n}^{i}.$
	
	Let $M=\oplus_{\lambda_{\in\mathbb{C}}}M_{\lambda}$ be a $V$-module.
	The restricted dual of $M$ is defined by $M'=\oplus_{\lambda_{\in\mathbb{C}}}M_{\lambda}^{\ast}$
	where $M_{\lambda}^{\ast}=\text{Hom}_{\mathbb{C}}\left(M_{\lambda},\mathbb{C}\right).$
	It was proved in \cite{FHL} that $M'=\left(M',Y_{M'}\right)$ is
	naturally a $V$-module such that
	\[
	\left\langle Y_{M'}\left(v,z\right)f,u\right\rangle =\left\langle f,Y_{M}\left(e^{zL\left(1\right)}\left(-z^{-2}\right)^{L\left(0\right)}v,z^{-1}\right)u\right\rangle ,
	\]
	for $v\in V,$ $f\in M'$ and $u\in M$, and $\left(M'\right)'\cong M$.
	Moreover, if $M$ is irreducible, so is $M'$. A $V$-module $M$
	is said to be \emph{self dual} if $M\cong M'$.
	
	\begin{definition} A vertex operator algebra $V$ is said to be \emph{$C_{2}$-cofinite}
		if $V/C_{2}(V)$ is finite dimensional, where $C_{2}(V)=\langle v_{-2}u|v,u\in V\rangle.$
	\end{definition}
	
	\begin{definition}A vertex operator algebra $V=\oplus_{n\in\mathbb{Z}}V_{n}$
		is said to be of \emph{CFT type} if $V_{n}=0$ for negative $n$ and
		$V_{0}=\mathbb{C}1$. \end{definition}
	
	\begin{definition} Let $\left(V,Y\right)$ be a vertex operator algebra
		and let $\left(M^{i},Y^{i}\right),\ \left(M^{j},Y^{j}\right)$ and
		$\left(M^{k},Y^{k}\right)$ be three $V$-modules. An \emph{intertwining
			operator of type $\left(\begin{array}{c}
				M^{k}\\
				M^{i}\ M^{j}
			\end{array}\right)$} is a linear map
		\begin{gather*}
			\mathcal{Y}\left(\cdot,z\right):\ M^{i}\to\text{\ensuremath{\mbox{Hom}\left(M^{j},\ M^{k}\right)\left\{ z\right\} }}\\
			u\mapsto\mathcal{Y}\left(u,z\right)=\sum_{n\in\mathbb{Q}}u_{n}z^{-n-1}
		\end{gather*}
		satisfying:
		
		(1) for any $u\in M^{i}$ and $v\in M^{j}$, $u_{n}v=0$ for $n$
		sufficiently large;
		
		(2) $\mathcal{Y}(L_{-1}v,\ z)=\left(\frac{d}{dz}\right)\mathcal{Y}\left(v,z\right)$
		for $v\in M^{i}$;
		
		(3) (Jacobi Identity) for any $u\in V,\ v\in M^{i}$,
		\begin{alignat*}{1}
			& z_{0}^{-1}\delta\left(\frac{z_{1}-z_{2}}{z_{0}}\right)Y^{k}\left(u,z_{1}\right)\mathcal{Y}\left(v,z_{2}\right)-z_{0}^{-1}\delta\left(\frac{-z_{2}+z_{1}}{z_{0}}\right)\mathcal{Y}\left(v,z_{2}\right)Y^{j}\left(u,z_{1}\right)\\
			& =z_{2}^{-1}\left(\frac{z_{1}-z_{0}}{z_{2}}\right)\mathcal{Y}\left(Y^{i}\left(u,z_{0}\right)v,z_{2}\right).
		\end{alignat*}
		The space of all intertwining operators of type $\left(\begin{array}{c}
			M^{k}\\
			M^{i}\ M^{j}
		\end{array}\right)$ is denoted $I_{V}\left(\begin{array}{c}
			M^{k}\\
			M^{i}\ M^{j}
		\end{array}\right)$. Without confusion, we also denote it by $I_{i,j}^{k}.$ Let $N_{i,\ j}^{k}=\dim I_{i,j}^{k}$.
		These integers $N_{i,j}^{k}$ are called the \emph{fusion rules}.
	\end{definition}
	
	% The following proposition was proved in \cite{ADL}:
	
	% \begin{proposition} \label{restriction of fusion rules} Let $V$
		% be a vertex operator algebra and let $W^{1}$, $W^{2}$, $W^{3}$
		% be $V$-modules among which $W^{1}$ and $W^{2}$ are irreducible.
		% Suppose that $V_{0}$ is a vertex operator subalgebra of $V$ (with
		% the same Virasoro element) and that $N^{1}$ and $N^{2}$ are irreducible
		% $V_{0}$-modules of $W^{1}$ and $W^{2}$, respectively. Then the
		% restriction map from $I_{V}\left(_{W^{1}\ W^{2}}^{\ \ W^{3}}\right)$
		% to $I_{V_{0}}\left(_{N^{1}\ N^{2}}^{\ \ W^{3}}\right)$ is injective.
		% In particular,
		% \[
		% \dim I_{V}\left(_{W^{1\ }W^{2}}^{\ \ W^{3}}\right)\le\dim I_{V_{0}}\left(_{N^{1}\ N^{2}}^{\ \ W^{3}}\right).
		% \]
		% \end{proposition}
	
	Let $V^{1}$ and $V^{2}$ be vertex operator algebras. Let $M^{i}$
	, $i=1,2,3$, be $V^{1}$-modules, and $N^{i}$, $i=1,2,3$, be $V^{2}$-modules.
	Then $M^{i}\otimes N^{i}$, $i=1,2,3$, are $V^{1}\otimes V^{2}$-modules
	by \cite{FHL}. The following property was given in \cite{ADL}:
	
	\begin{proposition} \label{fusion of tensor product}If $N_{M^{1},M^{2}}^{M^{3}}<\infty$
		or $N_{N^{1},N^{2}}^{N^{3}}<\infty,$ then
		\[
		N_{M^{1}\otimes N^{1},M^{2}\otimes N^{2}}^{M^{3}\otimes N^{3}}=N_{M^{1},M^{2}}^{M^{3}}N_{N^{1},N^{2}}^{N^{3}}.
		\]
	\end{proposition}
	
	Let $M^{1}$ and $M^{2}$ be $V$-modules. A fusion product for the
	ordered pair $\left(M^{1},M^{2}\right)$ is a pair $\left(M,F\left(\cdot,z\right)\right)$
	which consists of a $V$-module $M$ and an intertwining operator
	$F\left(\cdot,z\right)$ of type $\left(_{M^{1},M^{2}}^{\ \ M}\right)$
	such that the following universal property holds: For any $V$-module
	$W$ and any intertwining operator $I\left(\cdot,z\right)$ of type
	$\left(_{M^{1},M^{2}}^{\ \ W}\right)$, there exists a unique $V$-homomorphism
	$\phi$ from $M$ to $W$ such that $I\left(\cdot,z\right)=\phi\circ F\left(\cdot,z\right).$
	It is clear from the definition that if a tensor product of $M^{1}$
	and $M^{2}$ exists, it is unique up to isomorphism. In this case,
	we denote the \emph{fusion product} by $M^{1}\boxtimes_{V}M^{2}.$
	
	The basic result is that the fusion product exists if $V$ is rational.
	Let $M,N$ be irreducible $V$ -modules, we shall often consider the
	fusion product
	\[
	M\boxtimes_{V}N=\sum_{W}N_{M,\ N}^{W}W
	\]
	where $W$ runs over the set of equivalence classes of irreducible
	$V$-modules.
	
	\begin{definition} Let $V$ be a simple vertex operator algebra.
		A simple $V$-module $M$ is called \emph{simple current} if for
		any irreducible $V$-module $W$, $M\boxtimes_{V}W$ exists and is
		also a simple $V$-module. \end{definition}
	
	The following proposition is from \cite{FHL}:
	
	\begin{proposition} \label{extension property} Let $V$ be a vertex
		operator algebra and $V'$ be its restricted dual. For $u,v,w\in V$
		and $t\in V'$, we have the following equality of rational functions
		
		\begin{equation}
			\iota_{12}^{-1}\left\langle t,Y\left(u,z_{1}\right)Y\left(v,z_{2}\right)w\right\rangle =\iota_{21}^{-1}\left\langle t,Y\left(v,z_{2}\right)Y\left(u,z_{1}\right)w\right\rangle \label{communitivity}
		\end{equation}
		
		\begin{equation}
			\iota_{12}^{-1}\left\langle t,Y\left(u,z_{1}\right)Y\left(v,z_{2}\right)w\right\rangle =\left(\iota_{20}^{-1}\left\langle t,Y\left(Y\left(u,z_{0}\right)v,z_{2}\right)w\right\rangle \right)|_{z_{0}=z_{1}-z_{2}}\label{associavity}
		\end{equation}
		where $\iota_{12}^{-1}f\left(z_{1},z_{2}\right)$ denotes the formal
		power expansion of an analytic function $f\left(z_{1},z_{2}\right)$
		in the domain $\left|z_{1}\right|>\left|z_{2}\right|$ . \end{proposition}
	
	The following result about bilinear form on $V$ is from \cite{Li}:
	
	\begin{theorem}\label{bilinear form}The space of invariant bilinear
		forms on $V$ is isomorphic to the space
		\[
		\left(V_{0}/L\left(1\right)V_{1}\right)^{*}=\mbox{Hom}_{\mathbb{C}}\left(V_{0}/L\left(1\right)V_{1},\mathbb{C}\right).
		\]
	\end{theorem}
	
	\subsection{Quantum Galois theory}
	
	Now we recall quantum Galois theory and quantum dimensions from \cite{DM}
	and \cite{DJX}. Let $V$ be a simple vertex operator algebra and
	$G$ a finite and faithful group of automorphisms of $V$. Let $\text{Irr}\left(G\right)$
	be the set of simple characters $\chi$ of $G$. As $\mathbb{C}G$-module,
	each homogeneous space $V_{n}$ of $V$ is finite dimensional, and
	$V$ can be decomposed into a direct sum of graded subspaces
	\[
	V=\oplus_{\chi\in\text{Irr}\left(G\right)}V^{\chi},
	\]
	where $V^{\chi}$ is the subspace of $V$ on which $G$ acts according
	to the character $\chi$. The following theorem is from \cite{DM,DLM4}. 
	% also see \cite{DLM}.
	
	\begin{theorem} \label{classical galois theory}Suppose that $V$
		is a simple vertex operator algebra and that $G$ is a finite and
		faithful solvable group of automorphisms of $V$. Then the following
		hold:
		
		(i) Each $V^{\chi}$ is nonzero;
		
		(ii) For $\chi\in\text{Irr}\left(G\right)$, each $V^{\chi}$ is a
		simple module for the $G$-graded vertex operator algebra $\mathbb{C}G\otimes V^{G}$
		of the form
		
		\[
		V^{\chi}=M_{\chi}\otimes V_{\chi},
		\]
		where $M_{\chi}$ is the simple $\mathbb{C}G$-module affording $\chi$
		and where $V_{\chi}$ is a simple $V^{G}$-module.
		
		(iii) The map $M_{\chi}\mapsto V_{\chi}$ is a bijection from the
		set of simple $\mathbb{C}G$-modules to the set of (inequivalent)
		simple $V^{G}$-modules which are contained in $V$.
		
	\end{theorem}
	
	Now we recall the notion of quantum dimension from \cite{DJX}. Let
	$V$ be a vertex operator algebra of CFT type and $M$ a $V$-module,
	the formal character of $M=\oplus_{n\in\mathbb{Z}_{+}}M_{\lambda+n}$
	is defined to be
	\[
	Ch_{q}M=\mbox{tr}_{q}M=\text{tr}q^{L(0)-c/24}=q^{\lambda-c/24}\sum_{n\in\mathbb{Z}_{+}}(\dim M_{\lambda+n})q^{n}
	\]
	where $\lambda$ is the conformal weight of $M$. The quantum dimension
	of $M$ over $V$ is defined as:
	\[
	q\dim_{V}M=\lim_{q\to1^{-}}\frac{Ch_{q}M}{Ch_{q}V}.
	\]
	The following result is from Theorem 6.3 in \cite{DJX}:
	
	\begin{theorem} \label{quantum dimension and orbifold module}Let
		$V$ be a rational and $C_{2}$-cofinite simple vertex operator algebra.
		Assume $V$ is $g$-rational and the conformal weight of any irreducible
		$g$-twisted $V$-module is positive except for $V$ itself for all
		$g\in G$. Then
		\[
		q\dim_{V^{G}}V_{\chi}=\dim W_{\chi}.
		\]
	\end{theorem}
	
	For convenience, from now on, we say a vertex operator algebra $V$
	is ``good'' if it satisfies the following conditions: $V$ is a
	rational and $C_{2}$-cofinite simple vertex operator algebra of CFT
	type with $V\cong V'$. Let $M^{0},\ M^{1},\ \cdots,\ M^{d}$ be all
	the inequivalent irreducible $V$-modules with $M^{0}\cong V$ . The
	corresponding conformal weights $\lambda_{i}$ satisfy $\lambda_{i}>0$
	for $0<i\le d$.
	
	The following properties of quantum dimensions are from \cite{DJX}
	:
	
	\begin{proposition} \label{product property of quantum dimension }\label{quantum dimension and simple current}
		Let $V$ be a ``good'' vertex operator algebra. Then
		
		(i) $q\dim_{V}\left(M^{i}\boxtimes M^{j}\right)=q\dim_{V}M^{i}\cdot q\dim_{V}M^{j},$
		$\forall i,j$.
		
		(ii) A $V$-module $M^{i}$ is a simple current if and only if $q\dim_{V}M^{i}=1$.
		
		(iii) $q\dim_{V}M^{i}\in\left\{ 2\cos\left(\pi/n\right)|n\ge3\right\} \cup\left\{ a|2\le a<\infty,a\ is\ algebraic\right\} .$
		
	\end{proposition}
	
	\begin{definition} Let $V$ be a vertex operator algebra with finitely
		many inequivalent irreducible modules $M^{0},\cdots,M^{d}$. The \emph{global
			dimension} of $V$ is defined as
		\[
		\text{glob}\left(V\right)=\sum_{i=0}^{d}\left(q\dim_{V}M^{i}\right)^{2}.
		\]
	\end{definition}
	
	\begin{remark} \label{product property of global dimension}Let $U$
		and $V$ be ``good'' vertex operator algebras, $M$ be a $U$-module
		and $N$ be a $V$-module. Then Lemma 2.10 of \rm{\cite{ADJR}} gives
		
		\[
		q\dim_{U\otimes V}M\otimes N=q\dim_{U}M\cdot q\dim_{V}N,
		\]
		\[
		\text{glob}\left(U\otimes V\right)=\text{glob}\left(U\right)\cdot\text{glob}\left(V\right).
		\]
	\end{remark}
	
	Let $V$ be a vertex operator algebra, recall that a simple vertex
	operator algebra containing $V$ is called an \emph{extension} $U$
	of $V$. Now we have the following theorem \cite{ABD,HKL,ADJR}:
	
	\begin{theorem} \label{rationality of extesnion of VOA } \label{global dim of VOA and subVOA}
		Let $V$ be a ``good'' vertex operator algebra. Let $U$ be a simple
		vertex operator algebra which is an extension of $V$. Then $U$ is
		also ``good'' and
		\[
		\text{glob}\left(V\right)=\text{glob}\left(U\right)\cdot\left(q\dim_{V}\left(U\right)\right)^{2}.
		\]
		
	\end{theorem}
	
	\subsection{\label{subsec:The-unitary-series}The unitary series of the Virasoro
		VOAs}
	
	Now we recall notations about unitary minimal models of Virasoro algebra
	from \cite{FFK}. The models are parameterized by a complex number
	$\alpha_{-}^{2}$, related to the central charge of the Virasoro algebra
	by $c=13-6\alpha_{-}^{2}-6\alpha_{-}^{-2}$ where $\alpha_{-}^{2}=\frac{p}{p'}$
	and $\left|p-p'\right|=1.$ Without loss of generality, we write $p'=p+1$
	and denote $c_{p}=1-\frac{6}{p\left(p+1\right)}$ with $p=2,3,4,\cdots$.
	The label $I$ stands for a pair $\left(i',i\right)$ of positive
	integers and the corresponding highest weight is
	\begin{equation}
		h_{I}=h_{\left(i'\ i\right)}^{\left(p\right)}=\frac{1}{4}\left(i'^{2}-1\right)\alpha_{-}^{2}-\frac{1}{2}\left(i'i-1\right)+
		\frac{1}{4}\left(i^{2}-1\right)\alpha_{-}^{-2}=\frac{\left(pi'-\left(p+1\right)i\right)^{2}-1}{4p\left(p+1\right)}.\label{minimal model notation}
	\end{equation}
	for $1\le i'\le p,$ $1\le i\le p-1.$ We denote such unitary minimal
	models of Virasoro algebra by $L\left(c_{p},h_{\left(i',i\right)}^{\left(p\right)}\right).$
	
	\begin{remark} \label{pairs corr. to U's weights for unitary model}Use
		the above notation, we see that the central charge of the model $L\left(\frac{25}{28},0\right)$
		corresponds to the parameter $\alpha_{-}^{2}=\frac{7}{8}$ with $p=7,$
		$p'=8$. The highest weights for irreducible $L\left(\frac{25}{28},0\right)$-modules
		are
		
		\begin{equation}
			\left\{ 0,\frac{5}{32},\frac{3}{4},\frac{57}{32},\frac{13}{4},\frac{165}{32},\frac{15}{2},\frac{5}{14},\frac{3}{224},
			\frac{3}{28},\frac{143}{224},\frac{45}{28},\frac{675}{224},\frac{34}{7},\frac{9}{7},\frac{99}{224},\frac{1}{28},\frac{15}{224},\frac{25}{28},\frac{323}{224},
			\frac{39}{14}\right\} .\label{highest weights for 25/28}
		\end{equation}
		In particular, the pairs $\left(1,1\right)$, $\left(1,5\right)$
		and $\left(1,3\right)$ correspond to the highest weights $0$, $\frac{34}{7}$
		and $\frac{9}{7}$ respectively. \end{remark}
	
	Also note that the fusion rules for irreducible $L\left(c_{p},0\right)$-modules
	are as follows \cite{W}:
	
	\begin{definition} An ordered triple of pairs of integers $\left(\left(i',i\right),\left(j',j\right),\left(k',k\right)\right)$
		is called \emph{admissible }if $1\le i',j',k'\le p+1,1\le i,j,k\le p$,
		$i'+j'+k'<2\left(p+1\right),$ $i+j+k<2p$, $i'<j'+k'$, $j'<i'+k'$,
		$k'<i'+j'$, $i<j+k$, $j<i+k$, $k<i+j$, and the sums $i'+j'+k'$,
		$i+j+k$ are odd. \end{definition}
	
	\begin{proposition} \label{fusion rules of virasoro modules}The
		fusion rules between $L\left(c_{p},0\right)$-modules $L\left(c_{p},h_{\left(i',i\right)}^{\left(p\right)}\right),$
		$L\left(c_{p},h_{\left(j',j\right)}^{\left(p\right)}\right)$ are
		\[
		L\left(c_{p},h_{\left(i',i\right)}^{\left(p\right)}\right)\boxtimes L\left(c_{p},h_{\left(j',j\right)}^{\left(p\right)}\right)=
		\sum_{\left(k',k\right)}N_{\left(i',i\right),\left(j',j\right)}^{\left(k',k\right)}L\left(c_{p},h_{\left(k',k\right)}^{\left(p\right)}\right),
		\]
		where $N_{\left(i',i\right),\left(j',j\right)}^{\left(k',k\right)}$
		is $1$ iff $\left(\left(i',i\right),\left(j',j\right),\left(k',k\right)\right)$
		is an admissible triple of pairs and $0$ otherwise. \end{proposition}
	
	\subsection{\label{subsec:Braiding-matrix}Braiding matrices}
	
	Now we recall four point functions. Let $V$ be a rational and $C_{2}$-cofinite
	vertex operator algebra of CFT type and $V\cong V'$. Let $M^{a_{1}},M^{a_{2}},M^{a_{3}},M^{a_{4}}$
	be four irreducible $V$-modules. By Lemma 4.1 in \cite{H1} we know
	that for $u_{a_{i}}\in M^{a_{i}},$
	
	\[
	\left\langle u_{a_{4}'},\mathcal{Y}_{a_{1},a_{5}}^{a_{4}}\left(u_{a_{1}},z_{1}\right)\mathcal{Y}_{a_{2},a_{3}}^{a_{5}}
	\left(u_{a_{2}},z_{2}\right)u_{a_{3}}\right\rangle ,
	\]
	
	\[
	\left\langle u_{a_{4}'},\mathcal{Y}_{a_{2},a_{6}}^{a_{4}}\left(u_{a_{1}},z_{2}\right)\mathcal{Y}_{a_{1},a_{3}}^{a_{6}}\left(u_{a_{1}},z_{1}\right)u_{a_{3}}\right\rangle
	\]
	are analytic on $\left|z_{1}\right|>\left|z_{2}\right|>0$ and $\left|z_{2}\right|>\left|z_{1}\right|>0$
	respectively, and can both be analytically extended to multi-valued
	analytic functions on
	\[
	R=\left\{ \left(z_{1},z_{2}\right)\in\mathbb{C}^{2}|z_{1},z_{2}\not=0,z_{1}\not=z_{2}\right\} .
	\]
	One can lift the multi-valued functions on $R$ to single-valued functions
	on the universal covering $\tilde{R}$ to $R$ as in \cite{H1}. We
	use
	
	\[
	E\left\langle u_{a_{4}'},\mathcal{Y}_{a_{1},a_{5}}^{a_{4}}\left(u_{a_{1}},z_{1}\right)\mathcal{Y}_{a_{2},a_{3}}^{a_{5}}\left(u_{a_{2}},z_{2}\right)u_{a_{3}}\right\rangle
	\]
	and
	\[
	E\left\langle u_{a_{4}'},\mathcal{Y}_{a_{2},a_{6}}^{a_{4}}\left(u_{a_{1}},z_{2}\right)\mathcal{Y}_{a_{1},a_{3}}^{a_{6}}\left(u_{a_{1}},z_{1}\right)u_{a_{3}}\right\rangle
	\]
	to denote those analytic functions.
	
	Let $\left\{ \mathcal{Y}_{a,b;i}^{c}|i=1,\cdots,N_{a,b}^{c}\right\} $
	be a basis of $I_{a,b}^{c}$. From \cite{H2},
	
	\[
	\left\{ E\left\langle u_{a_{4}'},\mathcal{Y}_{a_{1},a_{5};i}^{a_{4}}\left(u_{a_{1}},z_{1}\right)\mathcal{Y}_{a_{2},a_{3};j}^{a_{5}}\left(u_{a_{2}},z_{2}\right)u_{a_{3}}\right\rangle |i=1,\cdots,N_{a_{1},a_{5}}^{a_{4}},j=1,\cdots,N_{a_{2},a_{3}}^{a_{5}},\forall a_{5}\right\}
	\]
	is a linearly independent set. Fix a basis of intertwining operators.
	It was proved in \cite{KZ,TK} that
	\begin{alignat*}{1}
		& span\left\{ E\left\langle u_{a_{4}'},\mathcal{Y}_{a_{3},\mu;i}^{a_{4}}\left(u_{a_{3}},z_{1}\right)\mathcal{Y}_{a_{2},a_{1};j}^{\mu}\left(u_{a_{2}},z_{2}\right)u_{a_{1}}\right\rangle |i,j,\mu\right\} \\
		& =span\left\{ E\left\langle u_{a_{4}'},\mathcal{Y}_{a_{2},\gamma;k}^{a_{4}}\left(u_{a_{2}},z_{2}\right)\mathcal{Y}_{a_{3},a_{1};l}^{\gamma}\left(u_{a_{3}},z_{1}\right)u_{a_{1}}\right\rangle |k,l,\gamma\right\} ,
	\end{alignat*}
	where $u_{a_{i}}\in M^{a_{i}}$. Then there exists $\left(B_{a_{4,}a_{1}}^{a_{3},a_{2}}\right)_{\mu,\gamma}^{i,j;k,l}\in\mathbb{C}$
	such that
	\begin{alignat}{1}
		& E\left\langle u_{a_{4}'},\mathcal{Y}_{a_{3},\mu;i}^{a_{4}}\left(u_{a_{3}},z_{1}\right)\mathcal{Y}_{a_{2},a_{1};j}^{\mu}\left(u_{a_{2}},z_{2}\right)u_{a_{1}}\right\rangle \nonumber \\
		& =\sum_{k,l,\gamma}\left(B_{_{a_{4},a_{1}}}^{a_{3,}a_{2}}\right)_{\mu,\gamma}^{i,j;k,l}E\left\langle u_{a_{4}'},\mathcal{Y}_{a_{2},\gamma;k}^{a_{4}}\left(u_{a_{2}},z_{1}\right)\mathcal{Y}_{a_{3},a_{1};l}^{\gamma}\left(u_{a_{3}},z_{2}\right)u_{a_{1}}\right\rangle \label{braiding matrix property}
	\end{alignat}
	(see \cite{H1}). $B_{_{a_{4},a_{1}}}^{a_{3,}a_{2}}$ is called
	the\emph{ braiding matrix}.
	
	% Let $b_{1},b_{2},b_{3},b_{4}$ be four irreducible $L\left(\frac{25}{28},0\right)$-modules.
	% Fix a basis $\left\{ \mathcal{\overline{Y}}_{a,b;i}^{c}|i=1,\cdots,N_{a,b}^{c}\right\} $
	% of intertwining operators of $I_{L\left(\frac{25}{28},0\right)}\left(_{L\left(\frac{25}{28},a\right),L\left(\frac{25}{28},b\right)}^{L\left(\frac{25}{28},c\right)}\right)$
	% with $N_{a,b}^{c}=\dim I_{L\left(\frac{25}{28},0\right)}\left(_{L\left(\frac{25}{28},a\right),L\left(\frac{25}{28},b\right)}^{L\left(\frac{25}{28},c\right)}\right)$
	% as in \cite{FFK}. Then there exists a matrix $\left(\tilde{B}_{_{b_{4},b_{1}}}^{b_{3},b_{2}}\right)_{\mu,\gamma}^{i,j;k,l}\in\mathbb{C}$
	% such that
	% \begin{alignat}{1}
		% & E\left\langle u_{b_{4}'},\mathcal{\overline{Y}}_{b_{3},\mu;i}^{b_{4}}\left(u_{b_{3}},z_{1}\right)\mathcal{\overline{Y}}_{b_{2},b_{1};j}^{\mu}\left(u_{b_{2}},z_{2}\right)u_{b_{1}}\right\rangle \nonumber \\
		% & =\left(\tilde{B}_{_{b_{4},b_{1}}}^{b_{3,}b_{2}}\right)_{\mu,\gamma}^{i,j;k,l}E\left\langle u_{b_{4}'},\mathcal{\overline{Y}}_{b_{2},\gamma;k}^{b_{4}}\left(u_{b_{2}},z_{2}\right)\mathcal{\overline{Y}}_{b_{3},b_{1};l}^{\gamma}\left(u_{b_{3}},z_{1}\right)u_{b_{1}}\right\rangle \label{tilde=00007BB=00007D}
		% \end{alignat}
	% by (\ref{braiding matrix property}).Now we recall some formulas about
	% minimal models of Virasoro vertex operator algebra given in \cite{FFK}.
	% We will use these formulas to prove some properties of braiding matrices,
	% which will be needed in the proof of uniqueness of the structure of
	% the vertex operator algebra $\mathcal{U}$.
	
	Recall that we have seen $\alpha_{-}^{2}=\frac{p}{p'}$ in Section
	\ref{subsec:The-unitary-series}. Now let $\alpha_{+}^{2}=-\frac{p'}{p}$,
	$x=\exp\left(2\pi i\alpha_{+}^{2}\right)$, $y=\exp\left(2\pi i\alpha_{-}^{2}\right)$,
	$\left[l\right]=x^{l/2}-x^{-l/2}$, $\left[l'\right]=y{}^{l'/2}-y{}^{-l'/2}.$
	Now we fix central charge $c_{p}$, denote $L\left(c_{p},h_{\left(i',i\right)}^{\left(p\right)}\right)$
	by $\left(i',i\right)$. Let $\left(a',a\right)$, $\left(m',m\right)$,
	$\left(n',n\right)$, $\left(c',c\right)$, $\left(b',b\right)$,
	$\left(d',d\right)$ be irreducible $L\left(c_{p},0\right)$-modules,
	the braiding matrices of screened vertex operators have the almost
	factorized form (cf. (2.19) of \cite{FFK}):
	\begin{alignat}{1}
		& \left(\tilde{B}_{\left(m',m\right),\left(n',n\right)}^{\left(a',a\right),\left(c',c\right)}\right){}_{\left(b',b\right),\left(d',d\right)}\nonumber \\
		& =i^{-\left(m'-1\right)\left(n-1\right)-\left(n'-1\right)\left(m-1\right)}\left(-1\right)^{1/2\left(a-b+c-d\right)\left(n'+m\right)+1/2\left(a'-b'+c'-d'\right)\left(n+m\right)}\label{FFK 2.19}\\
		& \cdot r\left(a',m',n',c'\right)_{b',d'}\cdot r\left(a,m,n,c\right)_{b,d},\nonumber
	\end{alignat}
	where the nonvanishing matrix elements of $r$-matrices are
	
	\begin{gather}
		r\left(a,1,n,c\right)_{a,c}=r\left(a,m,1,c\right)_{c,a}=1,\nonumber \\
		r\left(l\pm2,2,2,l\right)_{l\pm1,l\pm1}=x^{1/4},\nonumber \\
		r\left(l,2,2,l\right)_{l\pm1,l\pm1}=\mp x^{-1/4\mp l/2}\frac{\left[1\right]}{\left[l\right]},\nonumber \\
		r\left(l,2,2,l\right)_{l\pm1,l\mp1}=x^{-1/4}\frac{\left[l\pm1\right]}{\left[l\right]},\label{FFK 2.20}
	\end{gather}
	and the other $r$-matrices are given by the recursive relation
	\begin{gather}
		r\left(a,m+1,n,c\right)_{b,d}=\sum_{d_{1}\ge1}r\left(a,2,n,d_{1}\right)_{a_{1},d}\cdot r\left(a_{1},m,n,c\right)_{b,d_{1}},\nonumber \\
		r\left(a,m,n+1,c\right)_{b,d}=\sum_{d_{1}\ge1}r\left(a,m,2,c_{1}\right)_{b,d_{1}}\cdot r\left(d_{1},m,n,c\right)_{c_{1},d},\label{FFK 2.21}
	\end{gather}
	for any choice of $a_{1}$ and $c_{1}$ compatible with the fusion
	rules. The $r'$ matrices are given by the same formulas with the
	replacement $x\to x'$, $\left[\ \ \ \ \right]\to\left[\ \ \ \ \right]'.$
	
	\subsection{GKO construction of the unitary Virasoro VOA}
	
	Let $e,f$ and $h$ be the generators of $\mathfrak{sl}_{2}(\mathbb{C})$
	such that
	\[
	\left[e,f\right]=h,\ \left[h,e\right]=2e,\ \left[h,f\right]=-2f.
	\]
	Let $\left\langle \cdot,\cdot\right\rangle $ be the standard invariant
	bilinear form on $\mathfrak{sl}_{2}(\mathbb{C})$ defined by
	\[
	\left\langle h,h\right\rangle =2,\ \left\langle e,f\right\rangle =1,\ \left\langle e,e\right\rangle =\left\langle f,f\right\rangle =\left\langle h,e\right\rangle =\left\langle h,f\right\rangle =0.
	\]
	Let $\hat{\mathfrak{sl}}_{2}\left(\mathbb{C}\right)$ be the corresponding
	affine algebra of type $A_{1}^{(1)}$ and $\lambda_{0},\ \lambda_{1}$
	the fundamental weights for $\hat{\mathfrak{sl}}_{2}\left(\mathbb{C}\right)$.
	Denote
	\[
	\mathcal{L}(m,k)=\mathcal{L}\left(\left(m-k\right)\lambda_{0}+k\lambda_{1}\right)
	\]
	the irreducible highest weight module of $\hat{\mathfrak{sl}}_{2}\left(\mathbb{C}\right)$-module
	with highest weight $\left(m-k\right)\lambda_{0}+k\lambda_{1}$. It
	was proved in \cite{FZ} that $\mathcal{L}\left(m,0\right)$ has a
	natural vertex operator algebra structure for $m\in\mathbb{Z}_{+}$.
	The Virasoro vector $\omega^{m}$ of $\mathcal{L}\left(m,0\right)$
	is given by
	\[
	\omega^{m}=\frac{1}{2\left(m+2\right)}\left(\frac{1}{2}h_{-1}h+e_{-1}f+f_{-1}e\right)
	\]
	with central charge $\frac{3m}{m+2}$. Let $m\in\mathbb{Z}_{+}$,
	then $\mathcal{L}\left(m,0\right)$ is a rational vertex operator
	algebra and
	$$\left\{ \mathcal{L}\left(m,\ k\right)|k=0,1,\cdots,m\right\}$$
	is the set of all the irreducible 
	$\mathcal{L}\left(m,0\right)$-modules.
	Moreover, the fusion rules are given by
	
	\[
	\mathcal{L}\left(m,j\right)\boxtimes\mathcal{L}\left(m,k\right)=\sum_{i=\max\left\{ 0,j+k-m\right\} }^{\min\left\{ j,k\right\} }\mathcal{L}\left(m,j+k-2i\right).
	\]
	
	Let $\mathcal{L}\left(m,0\right)_{1}$ be the weight 1 subspace of
	$\mathcal{L}(m,0)$. Then $\mathcal{L}(m,0)_{1}$ has a structure
	of Lie algebra isomorphic to $\mathfrak{sl}_{2}(\mathbb{C})$ under
	$[a,b]=a_{0}b$, $\forall a,b\in\mathcal{L}(m,0)$. Let $h^{m}$,
	$e^{m}$, $f^{m}$ be the generators of $\mathfrak{sl}_{2}(\mathbb{C})$
	in $\mathcal{L}(m,0)_{1}$. Then $h^{m+1}:=h^{1}\otimes1+1\otimes h^{m},$
	$e^{m+1}:=e^{1}\otimes1+1\otimes e^{m}$ and $f^{m+1}:=f^{1}\otimes1+1\otimes f^{m}$
	generate a vertex operator subalgebra isomorphic to $\mathcal{L}\left(m+1,0\right)$
	in $\mathcal{L}\left(m,1\right)\otimes\mathcal{L}\left(m,0\right)$.
	It was proved in \cite{DL} and \cite{KR} that $\Omega^{m}:=\omega^{1}\otimes1+1\otimes\omega^{m}-\omega^{m+1}$
	also gives a Virasoro vector with central charge $c_{m+2}=1-6/\left(m+2\right)\left(m+3\right)$.
	Furthermore, $\omega^{m+1}$ and $\Omega^{m}$ are mutually commutative
	and $\Omega^{m}$ generates a simple Virasoro vertex operator algebra
	$L\left(c_{m+2},0\right)$. Therefore $\mathcal{L}\left(1,0\right)\otimes\mathcal{L}\left(m,0\right)$
	contains a vertex operator subalgebra isomorphic to $L\left(c_{m+2},0\right)\otimes\mathcal{L}\left(m+1,0\right)$.
	Note that both $L\left(c_{m+2},0\right)$ and $\mathcal{L}\left(m+1,0\right)$
	are rational and every $\mathcal{L}\left(1,0\right)\otimes\mathcal{L}\left(m,0\right)$-module
	can be decomposed into irreducible $L(c_{m+2},0)\otimes\mathcal{L}(m+1,0)$-submodules.
	We have the following decomposition \cite{GKO}:
	
	\begin{equation}
		\mathcal{L}\left(1,\epsilon\right)\otimes\mathcal{L}\left(m,n\right)=\bigoplus_{0\le s\le m+3,s\equiv n+\epsilon\ \rm{mod} 2}L\left(c_{m+2},h_{\left(s+1,n+1\right)}^{(m+2)}\right)\otimes\mathcal{L}\left(m+1,s\right)\label{eq:(2.3)}
	\end{equation}
	where $\epsilon=0,1$ and $0\le n\le m$. This is the GKO-construction
	of the unitary Virasoro vertex operator algebras.
	
	\subsection{Structure of 6A-algebra}
	%This part we give the definition of 6A-Algebra $\mathcal{U}_{\rm{6A}}$ and state results from \cite{DJY}. 
	%Their work on $\mathcal{U}_{\rm{6A}}$ makes our induction possible.
	$6A$-algebra $\mathcal{U}_{6A}$ is a
	vertex operator algebra of the Moonshine type generated by two Ising
	vectors whose inner product is $5/2^{10}.$
	
	A vector $v\in V_{2}$ is called a \emph{Virasoro
		vector with central charge $c_{v}$ }if it satisfies\emph{ $v_{1}v=2v$
	}and $v_{3}v=\frac{c_{v}}{2}\mathbf{1}$. Then the operators $L_{n}^{v}:=v_{n+1},\ n\in\mathbb{Z}$,
	satisfy the Virasoro commutation relation
	\[
	\left[L_{m}^{v},\ L_{n}^{v}\right]=\left(m-n\right)L_{m+n}^{v}+\delta_{m+n,\ 0}\frac{m^{3}-m}{12}c_{v}
	\]
	for $m,\ n\in\mathbb{Z}.$ A Virasoro vector $v\in V_{2}$ with central
	charge $1/2$ is called an \emph{Ising vector }if $v$ generates the
	Virasoro vertex operator algebra $L(1/2,\ 0)$.
	
	It is shown in \cite{C} that the structure of the subalgebra generated
	by two Ising vectors $e$ and $f$ in the algebra $V_{2}^{\natural}$
	depends on only the conjugacy class of $\tau_{e}\tau_{f}$, and the
	inner product $\left\langle e,f\right\rangle $ is given by the following
	table:
	\begin{center}
		\begin{tabular}{|c|c|c|c|c|c|c|c|c|c|}
			\hline
			$\left\langle \tau_{e}\tau_{f}\right\rangle ^{\mathbb{M}}$  & $1A$  & $2A$  & $3A$  & $4A$  & $5A$  & $6A$  & $3C$  & $4B$  & $2B$\tabularnewline
			\hline
			$\left\langle e,f\right\rangle $  & $1/4$  & $1/2^{5}$  & $13/2^{10}$  & $1/2^{7}$  & $3/2^{9}$  & $5/2^{10}$  & $1/2^{8}$  & $1/2^{8}$  & $0$\tabularnewline
			\hline
		\end{tabular}
		\par\end{center}
	
	
	Certain coset subalgebra of $V_{\sqrt{2}E_{8}}$ associated with extended
	$E_{8}$ diagram is constructed in \cite{LYY} by removing one node
	from the diagram. In each case, the coset subalgebra contains some
	Ising vectors and the coset subalgebra is generated by two Ising vectors
	with inner product the same as the number given in the table above.
	In particular, the coset subalgebra $\mathcal{U}_{6A}$ corresponding
	to the $6A$ case was constructed, i.e., the case with inner product
	$\frac{5}{2^{10}}.$ Let $\mathcal{V}$ be the $3A$-algebra, that
	is, the vertex operator algebra generated by two Ising vectors whose
	$\tau$-involutions generate $S_{3}$ and with inner product $\frac{13}{2^{10}}$.
	The candidates for $\mathcal{V}$ were given \cite{M} and it was
	proved in \cite{SY} that only one of these candidates actually exists
	and that there is unique vertex operator algebra structure on it.
	
	
	Now we recall the following results about the $3A$-algebra $\mathcal{V}$
	from \cite{SY} .
	
	\begin{lemma} \label{rationality of 3A}The $3A$-algebra $\mathcal{V}$
		is rational. \end{lemma}
	
	\begin{lemma} \label{U_3A modules} All the irreducible $\mathcal{V}$-modules
		are as follows:
		
		\begin{align*}
			\mathcal{V}=\mathcal{V}\left(0\right) & =\left(\left(L\left(\frac{4}{5},0\right)\oplus L\left(\frac{4}{5},3\right)\right)\otimes\left(L\left(\frac{6}{7},0\right)\oplus L\left(\frac{6}{7},5\right)\right)\right)\\
			& \oplus L\left(\frac{4}{5},\frac{2}{3}\right)^{+}\otimes L\left(\frac{6}{7},\frac{4}{3}\right)^{+}\oplus L\left(\frac{4}{5},\frac{2}{3}\right)^{-}\otimes L\left(\frac{6}{7},\frac{4}{3}\right)^{-},
		\end{align*}
		
		\begin{align*}
			\mathcal{V}\left(\frac{1}{7}\right) & =\left(L\left(\frac{4}{5},0\right)\oplus L\left(\frac{4}{5},3\right)\right)\otimes\left((L\left(\frac{6}{7},\frac{1}{7}\right)\oplus L\left(\frac{6}{7},\frac{22}{7}\right)\right)\\
			& \oplus L(\left(\frac{4}{5},\frac{2}{3}\right)^{+}\otimes L\left(\frac{6}{7},\frac{10}{21}\right)^{+}\oplus L\left(\frac{4}{5},\frac{2}{3}\right)^{-}\otimes L\left(\frac{6}{7},\frac{10}{21}\right)^{-},
		\end{align*}
		
		\begin{align*}
			\mathcal{V}\left(\frac{5}{7}\right) & =\left(L\left(\frac{4}{5},0\right)\oplus L\left(\frac{4}{5},3\right)\right)\otimes\left(L\left(\frac{6}{7},\frac{5}{7}\right)\oplus L\left(\frac{6}{7},\frac{12}{7}\right)\right)\\
			& \oplus L(\left(\frac{4}{5},\frac{2}{3}\right)^{+}\otimes L\left(\frac{6}{7},\frac{1}{21}\right)^{+}\oplus L\left(\frac{4}{5},\frac{2}{3}\right)^{-}\otimes L\left(\frac{6}{7},\frac{1}{21}\right)^{-},
		\end{align*}
		
		\begin{align*}
			\mathcal{V}\left(\frac{2}{5}\right) & =\left(L\left(\frac{4}{5},\frac{2}{5}\right)\oplus L\left(\frac{4}{5},\frac{7}{5}\right)\right)\otimes\left(L\left(\frac{6}{7},0\right)\oplus L\left(\frac{6}{7},5\right)\right)\\
			& \oplus L\left(\frac{4}{5},\frac{1}{15}\right)^{+}\otimes L(\left(\frac{6}{7},\frac{4}{3}\right)^{+}\oplus L\left(\frac{4}{5},\frac{1}{15}\right)^{-}\otimes L\left(\frac{6}{7},\frac{4}{3}\right)^{-},
		\end{align*}
		
		\begin{align*}
			\mathcal{V}\left(\frac{19}{35}\right) & =\left(L\left(\frac{4}{5},\frac{2}{5}\right)\oplus L\left(\frac{4}{5},\frac{7}{5}\right)\right)\otimes\left(L\left(\frac{6}{7},\frac{1}{7}\right)\oplus L\left(\frac{6}{7},\frac{22}{7}\right)\right)\\
			& \oplus L(\left(\frac{4}{5},\frac{1}{15}\right)^{+}\otimes L(\left(\frac{6}{7},\frac{10}{21}\right)^{+}\oplus L(\left(\frac{4}{5},\frac{1}{15}\right)^{-}\otimes L(\left(\frac{6}{7},\frac{10}{21}\right)^{-},
		\end{align*}
		\begin{align*}
			\mathcal{V}\left(\frac{39}{35}\right) & =\left(L\left(\frac{4}{5},\frac{2}{5}\right)\oplus L\left(\frac{4}{5},\frac{7}{5}\right)\right)\otimes\left(L\left(\frac{6}{7},\frac{5}{7}\right)\oplus L\left(\frac{6}{7},\frac{12}{7}\right)\right)\\
			& \oplus L(\left(\frac{4}{5},\frac{1}{15}\right)^{+}\otimes L(\left(\frac{6}{7},\frac{1}{21}\right)^{+}\oplus L(\left(\frac{4}{5},\frac{1}{15}\right)^{-}\otimes L\left(\frac{6}{7},\frac{1}{21}\right)^{-}.
		\end{align*}
		
	\end{lemma}
	
	
	It was proved in \cite{LYY} that $\mathcal{V}\subset\mathcal{U}$
	and as a module of $\mathcal{V}\otimes L\left(\frac{25}{28},0\right)$,
	
	\[
	\mathcal{U}_{6A}\cong\mathcal{V}\otimes L\left(\frac{25}{28},0\right)\oplus\mathcal{V}\left(\frac{1}{7}\right)\otimes L\left(\frac{25}{28},\frac{34}{7}\right)\oplus\mathcal{V}\left(\frac{5}{7}\right)\otimes L\left(\frac{25}{28},\frac{9}{7}\right).
	\]
	From here forward, we denote
	\begin{gather}
		P_{1}=\mathcal{V},\ P_{2}=\mathcal{V}\left(\frac{1}{7}\right),\ P_{3}=\mathcal{V}\left(\frac{5}{7}\right),\nonumber \\
		Q_{1}=L\left(\frac{25}{28},0\right),\ Q_{2}=L\left(\frac{25}{28},\frac{34}{7}\right),\ Q_{3}=L\left(\frac{25}{28},\frac{9}{7}\right),\label{PQ notation}
	\end{gather}
	and $U^{i}=P_{i}\otimes Q_{i}$, $i=1,2,3$. Then
	\[
	\mathcal{U}_{6A}\cong P_{1}\otimes Q_{1}\oplus P_{2}\otimes Q_{2}\oplus P_{3}\otimes Q_{3}=U^{1}\oplus U^{2}\oplus U^{3}.
	\]
	
	\begin{remark} \label{rationality-1} Since $\mathcal{V}\otimes L\left(\frac{25}{28},0\right)$
		is a rational and $C_{2}$-cofinite vertex operator algebra, it is
		straightforward to see that $\mathcal{U}$ is also rational and $C_{2}$-cofinite
		by \rm{\cite{ABD,HKL}}. \end{remark}
	
	\begin{remark} \label{self-dual}
		Since $\mathcal{U}_{1}=0$ and $\dim\mathcal{U}_{0}=1$ by Theorem
		\ref{bilinear form}, there is a unique bilinear form on $\mathcal{U}$
		and thus $\mathcal{U}'\cong\mathcal{U}$. Without loss of generality,
		we can identify $\mathcal{U}$ with $\mathcal{U}'$.
	\end{remark}   
	
	Now we recall the following results in \cite{DJY}.
	
	% \begin{lemma}\label{new structure} Let $(V,Y)$ be a vertex operator
		% algebra and $f:V\to V$ be a linear isomorphism such that $\sigma\left(1\right)=1,\sigma\left(\omega\right)=\omega$.
		% Then $(V,Y^{\sigma})$ is a vertex operator where
		% \[
		% Y^{\sigma}(u,z)=\sigma Y(\sigma^{-1}u,z)\sigma^{-1}
		% \]
		% and $(V,Y)\cong(V,Y^{\sigma})$. 
		% \end{lemma}
	
	\begin{theorem}\label{uniqueness of 6A-algebra}
		The vertex operator algebra structure on $6A$-algebra $\mathcal{U}_{6A}$
		over $\mathbb{C}$ is unique.     
	\end{theorem}
	
	For simplicity, we shall use $[h_{1},\ h_{2}]$
	to denote the module $\mathcal{V}(h_{1})\otimes L(\frac{25}{28},\ h_{2})$ in this section.
	
	\begin{lemma} \label{realization of VOA}We have the following decomposition:
		\begin{align*}
			& \mathcal{L}\left(3,0\right)\otimes\mathcal{L}\left(1,0\right)\otimes\mathcal{L}\left(1,0\right)\otimes\mathcal{L}\left(1,0\right)\oplus\mathcal{L}\left(3,3\right)\otimes\mathcal{L}\left(1,1\right)\otimes\mathcal{L}\left(1,0\right)\otimes\mathcal{L}\left(1,0\right)\\
			& \cong\left\{ \left[0,\ 0\right]\oplus\left[\frac{1}{7},\frac{34}{7}\right]\oplus\left[\frac{5}{7},\frac{9}{7}\right]\right\} \otimes\mathcal{L}\left(6,0\right)\\
			& \oplus\left\{ \left[0,\frac{3}{4}\right]\oplus\left[\frac{5}{7},\frac{1}{28}\right]\oplus\left[\frac{1}{7},\frac{45}{28}\right]\right\} \otimes\mathcal{L}\left(6,2\right)\\
			& \oplus\left\{ \left[0,\frac{13}{4}\right]\oplus\left[\frac{1}{7},\frac{3}{28}\right]\oplus\left[\frac{5}{7},\frac{15}{28}\right]\right\} \otimes\mathcal{L}\left(6,4\right)\\
			& \oplus\left\{ \left[0,\frac{15}{2}\right]\oplus\left[\frac{1}{7},\frac{5}{14}\right]\oplus\left[\frac{5}{7},\frac{39}{14}\right]\right\} \otimes\mathcal{L}\left(6,6\right)
		\end{align*}
		Thus $\text{\ensuremath{\mathcal{L}}}\left(3,0\right)\otimes\mathcal{L}\left(1,0\right)\otimes\mathcal{L}\left(1,0\right)\otimes\mathcal{L}\left(1,0\right)\oplus\mathcal{L}\left(3,3\right)\otimes\mathcal{L}\left(1,1\right)\otimes\mathcal{L}\left(1,0\right)\otimes\mathcal{L}\left(1,0\right)$
		and $V_{L}$ contain a vertex operator subalgebra isomorphic to
		\[
		\left[0,\ 0\right]\oplus\left[\frac{1}{7},\frac{34}{7}\right]\oplus\left[\frac{5}{7},\frac{9}{7}\right]
		\]
		which is isomorphic to $\mathcal{U}_{6A}$ from the uniqueness of $\mathcal{U}_{6A}$.
	\end{lemma} 
	
	\begin{lemma} \label{Real modules} The following list give 14 irreducible
		$\mathcal{U}_{6A}$-module.
		\begin{align*}
			M^{0} & =\left[0,0\right]\oplus\left[\frac{1}{7},\frac{34}{7}\right]\oplus\left[\frac{5}{7},\frac{9}{7}\right], & M^{1} & =\left[0,\frac{3}{4}\right]\oplus\left[\frac{1}{7},\frac{45}{28}\right]\oplus\left[\frac{5}{7},\frac{1}{28}\right],\\
			M^{2} & =\left[0,\frac{13}{4}\right]\oplus\left[\frac{1}{7},\frac{3}{28}\right]\oplus\left[\frac{5}{7},\frac{15}{28}\right], & M^{3} & =\left[0,\frac{15}{2}\right]\oplus\left[\frac{1}{7},\frac{5}{14}\right]\oplus\left[\frac{5}{7},\frac{39}{14}\right],\\
			M^{4} & =\left[0,\frac{165}{32}\right]\oplus\left[\frac{1}{7},\frac{3}{224}\right]\oplus\left[\frac{5}{7},\frac{323}{224}\right], & M^{5} & =\left[0,\frac{5}{32}\right]\oplus\left[\frac{1}{7},\frac{675}{224}\right]\oplus\left[\frac{5}{7},\frac{99}{224}\right],\\
			M^{6} & =\left[0,\frac{57}{32}\right]\oplus\left[\frac{1}{7},\frac{143}{224}\right]\oplus\left[\frac{5}{7},\frac{15}{224}\right], & M^{7} & =\left[\frac{2}{5},0\right]\oplus\left[\frac{19}{35},\frac{34}{7}\right]\oplus\left[\frac{39}{35},\frac{9}{7}\right],\\
			M^{8} & =\left[\frac{2}{5},\frac{3}{4}\right]\oplus\left[\frac{19}{35},\frac{45}{28}\right]\oplus\left[\frac{39}{35},\frac{1}{28}\right], & M^{9} & =\left[\frac{2}{5},\frac{13}{4}\right]\oplus\left[\frac{19}{35},\frac{3}{28}\right]\oplus\left[\frac{39}{35},\frac{15}{28}\right],\\
			M^{10} & =\left[\frac{2}{5},\frac{15}{2}\right]\oplus\left[\frac{19}{35},\frac{5}{14}\right]\oplus\left[\frac{39}{35},\frac{39}{14}\right], & M^{11} & =\left[\frac{2}{5},\frac{5}{32}\right]\oplus\left[\frac{19}{35},\frac{675}{224}\right]\oplus\left[\frac{39}{35},\frac{99}{224}\right],\\
			M^{12} & =\left[\frac{2}{5},\frac{57}{32}\right]\oplus\left[\frac{19}{35},\frac{143}{224}\right]\oplus\left[\frac{39}{35},\frac{15}{224}\right], & M^{13} & =\left[\frac{2}{5},\frac{165}{32}\right]\oplus\left[\frac{19}{35},\frac{3}{224}\right]\oplus\left[\frac{39}{35},\frac{323}{224}\right].
		\end{align*}
	\end{lemma}
	
	% \subsection{Tensor categories}
	% In this subsection, we recall some notions and facts in the theory of tensor categories from \cite{BK}, \cite{KO} and \cite{O}.
	
	% Let $\cc$ %$(\cc, \1_{\cc}, \otimes )$
	% be a monoidal category defined as in \cite{BK}  and  denote the unit object in $\cc$ by $\1_{\cc}$.
	% %\begin{definition}{\rm
	% For  an object $U$ in  $\mathcal{C}$, %and $U$ an object in $\mathcal{C}$,
	% a {\em (right) dual} of $U$ is an object $U^\ast$ with two morphisms
	% $$e_{U}: U^{\ast}\otimes U\longrightarrow \1_{\cc},$$
	% $$i_{U}: \1_{\cc} 
	% \longrightarrow  U\otimes U^{\ast},$$ 
	% such that the composition 
	% $$U\stackrel{i_U\otimes \id_U}{\longrightarrow} U\otimes U^{\ast}\otimes U
	% \stackrel{\id_U\otimes e_U}{\longrightarrow} U$$ is equal to $\id_{U}$ and the composition $$U^{\ast}\stackrel{\id_{U^{\ast}}\otimes i_U}{\longrightarrow} U^{\ast}\otimes U\otimes U^{\ast}\stackrel{e_U\otimes 
		% \id_{U^{\ast}}}{\longrightarrow} U^{\ast}$$ is equal to $\id_{U^{\ast}}$.
	% %}\end{definition}
	% Similarly, for any object $U\in \cc$ we could define the {\em left dual} $\,^*U$ of $U$.
	% If the dual of $U$ exist, then   the dual of $U$ is unique (see \cite{BK}). 
	% %Furthermore, we have the following results which will be used frequently in this paper (see \cite{BK}).
	% % Similarly, for any object $U\in \cc$ we could define the {\em left dual} $\,^*U$ of $U$.
	% \begin{proposition}\label{kc3}
% Let $\cc$ %$(\cc, \1_{\cc}, \otimes )$
% be a monoidal category and $V$ be an object in $\cc$. If $V$ has a dual $V^*$, then
% $$\begin{aligned}
	% &\Hom(U\ot V, W)=\Hom(U, W\ot V^*),\\
	% &\Hom(U, V\ot W)=\Hom(V^*\ot U,  W).
	% \end{aligned}
% $$
% \end{proposition}
% Recall \cite{BK} that a monoidal category $\mathcal{C}$
% is called {\em a braided tensor category} if there is a natural bifunctor isomorphism $c_{X, Y}: X\otimes Y\to Y\otimes X$ called
% a {\em braiding} subject to the hexagon axioms.
% A definition of modular tensor category can be found in \cite{BK}.

% This result comes from \cite{KO} and \cite{O}.
% \begin{theorem}\label{kc2}
% %(i) $\cf_A$ is a monoidal functor.\\
% %$\cf$ and $\cg$ are adjoint, i.e. $$\Hom_A(A\otimes M, N)=\Hom_{\cc}(M, \cg(N))$$ for any $M\in \cc$ and $N\in \Rep A$.
% Let $\cc$ be a modular tensor category and  $A$ be an algebra in $\mathcal{C}$. For any object $U\in \cc$ and left  $A$-module $M$, we have
% $$\Hom_{A}(A\ot U, M)=\Hom(U, M),$$
% where we denote $\Hom_{_{A}\cc}$  by $\Hom_A$. Similarly, for any object $U\in \cc$ and  $A$-module $M$, we have
% $$\Hom_{\cc_A}(U\ot A, M)=\Hom(U, M).$$
% \end{theorem}

\subsection{Mirror extensions for rational vertex operator algebras}
A vertex operator algebra $V$ is said to be {\em regular} if any weak $V$-module
$M$ is a direct sum of irreducible $V$-modules.
%}\end{definition}

We now assume that $V$ is a  vertex operator algebra satisfying the following conditions:

(1) $V$ is simple CFT type vertex operator algebra  and is self dual;% i.e. $V$ is rational and $C_2$-cofinite;

(2) $V$ is regular.

The following results were obtained in \cite{ABD}, \cite{L}.
\begin{theorem}
Let $V$ be a $CFT$ type vertex operator algebra. Then $V$ is regular if and only if $V$ is rational and $C_2$-cofinite.
\end{theorem}

We have the
following results which were proved in \cite{FHL}, \cite{DMZ} and \cite{ABD}.
\begin{theorem}\label{kvoa3}
Let $V^1, \cdots, V^p$ be  vertex operator algebras. If $V^1, \cdots, V^p$ are rational, then
$V^1\otimes \cdots \otimes V^p$  is rational and any irreducible $V^1\otimes \cdots \otimes V^p$-module is a tensor
product $M^1\otimes \cdots \otimes M^p$ , where $M^1, \cdots, M^p$ are
some irreducible modules for the vertex operator algebras $V^1,
\cdots, V^p,$ respectively. If $V^1, \cdots, V^p$ are $C_2$-cofinite, then $V^1\otimes \cdots \otimes V^p$  is $C_2$-cofinite.
In particular, if $V^1, \cdots, V^p$ are  vertex operator algebras satisfying conditions $(1)$ and $(2)$, then $V^1\otimes \cdots \otimes V^p$ satisfies conditions $(1)$ and $(2)$.
\end{theorem}    

% A {\em  $V$-module} is a weak $V$-module $M$ which carries a $\mathbb{C}$-grading induced by the spectrum of $L(0)$, that is  $M=\bigoplus_{\lambda\in\mathbb{C}}
% M_{\lambda}$ where
% $M_\lambda=\{w\in M|L(0)w=\lambda w\}$. Moreover one requires that $M_\lambda$ is
% finite dimensional and for fixed $\lambda\in\mathbb{C}$, $M_{\lambda+n}=0$
% for sufficiently small integer $n$.
% %\begin{definition}{\rm
% Similarly, a $V$-module $M$ is said to be
% \textit{irreducible} if $M$ has no non-trivial
% $V$-submodule and 

We use the following theorem in \cite{Lin} to prove planty of essential results.

Let $V^1$, $V^2$ be  vertex operator algebras satisfying conditions (1) and (2). Denote  the module categories of $V^1$, $V^2$ by  $\mathcal{C}_{V^1}$, $\mathcal{C}_{V^2}$, respectively. 
If $V$ is an extension vertex operator algebra of $V^1\otimes V^2$, it is known that $V^1\otimes V^2$ is rational by \cite{FHL}, 
\cite{DMZ} and \cite{ABD},  then  $V$ has the following decomposition as $V^1\otimes V^2$-module:
$$V=\oplus_{i\in I, j\in J}Z_{i,j}M^i\otimes N^j,$$
where $Z_{i, j}$, $(i\in I, j\in J)$,  are nonnegative integers  and $\{M^i|i\in I\}$ (resp. $\{N^j|j\in J\}$) are inequivalent 
irreducible $V^1$-modules (resp. $V^2$-modules) such that $Z_{i, j}\neq 0$  for some $(i,j)\in I\times J$. 
%Note that  $M^i$ (resp. $N^j$) may be isomorphic to another $M^{i_1}$ (resp. $N^{j_1}$).
In the following we  assume that $V$ is an  extension vertex operator algebra  of $V^1\otimes V^2$ such that $V$ is simple, self dual
%(E2) $Z_{i, j}= 1$ for any $i\in I, j\in J$,\\
and $\hom_{V^1\otimes V^2}(V^1\otimes N^j, V)=\mathbb{C}$ (resp. $\hom_{V^1\otimes V^2}(M^i\otimes V^2, V)=\mathbb{C}$) if and only if $N^j=V^2$ (resp. $M^i=V^1$).
%and $\Hom_{V^1\otimes V^2}(M^i\otimes V^2, V)=\C$ if and only if $M^i=V^1$.\\
Then we have

\begin{theorem}\label{c1}
(i) $Z_{i,j}=1$ if $Z_{i,j}\neq 0$. Moreover, this induces a bijection $\tau$ from $I$ to $J$ such that $\tau(i)=j$ if and only if $Z_{i, j}=1$. \\
(ii) The set $\{M^i| i\in I\}$ $($resp. $\{N^j| j\in J\})$ is closed under the tensor product of category $\mathcal{C}_{V^1}$ $($resp. $\mathcal{C}_{V^2})$. \\
%In particular, $\{M^i| i\in I\}$ $($resp. $\{N^j| j\in J\})$ form a braided monoidal subcategory $\cc_{V^1}^0$ $($resp. $\cc_{V^2}^0)$ of  $\cc_{V^1}$ $($resp. $\cc_{V^2})$.\\
(iii) For any $i_1, i_2,i_3\in I$,  $N_{M^{i_1}, M^{i_2}}^{M^{i_3}}=N_{(N^{\tau(i_1)})', (N^{\tau(i_2)})'}^{(N^{\tau(i_3)})'}$.
\end{theorem}

\section{Construction of $U_k^i$ and $\{\mathcal{U}_k\}$}
Let $\mathcal{U}_{0}=\mathcal{V}$ be the $3A$-algebra and $\mathcal{U}_1=\mathcal{U}_{\operatorname{6A}}$.
By GKO-construction and Lemma (\ref{realization of VOA}) we can write
$$\left(\mathcal{L}\left(3,0\right)\otimes\mathcal{L}\left(1,0\right)\otimes\mathcal{L}\left(1,0\right)\otimes\mathcal{L}\oplus\mathcal{L}\left(3,3\right)\otimes\mathcal{L}\left(1,1\right)\otimes\mathcal{L}\left(1,0\right)\right)\otimes\mathcal{L}(1,0)\cong\bigoplus_{i=0,\rm{even}}^6\mathcal{U}_{1,i}\otimes \mathcal{L}(6,i),$$
where $\mathcal{U}_{1,i}$ is defined by the summand before $\mathcal{L}(6,i)$ 
and, especially, $\mathcal{U}_{1,0}=\mathcal{U}_{1}$ is the commutant of $\mathcal{L}(6,0)$.

\begin{example}
% $M_{0}^0=\mathcal{V}$, $M_{0}^1=\mathcal{V}(\frac{5}{7})$,  $M_{0}^2=\mathcal{V}(\frac{1}{7})$ which are irreducible modules of $\mathcal{U}_0=\mathcal{V}$.    
It is shown in Lemma (\ref{realization of VOA}) that 
$$\mathcal{U}_1=\mathcal{U}_{1,0}= U_0^1\oplus U_0^3\oplus U_0^5
%\bigoplus_{i=0}^2[i,h_{(1,i)}^{(8)}]_0
=\mathcal{V}\otimes L(\frac{25}{28},0)\oplus\mathcal{V}(\frac{5}{7})\otimes L(\frac{25}{28},\frac{9}{7})\oplus \mathcal{V}(\frac{1}{7})\otimes L(\frac{25}{28},\frac{34}{7}),$$
% where $h_j$ is defined by the corresponding $h$ appear in the summand of $\mathcal{L}(6,0)$ in GKO-construction.
and
$$
\begin{aligned}
\mathcal{U}_{1,2}&=\left[0,\frac{3}{4}\right]\oplus\left[\frac{5}{7},\frac{1}{28}\right]\oplus\left[\frac{1}{7},\frac{45}{28}\right],\\
\mathcal{U}_{1,4}&=\left[0,\frac{13}{4}\right]\oplus\left[\frac{1}{7},\frac{3}{28}\right]\oplus\left[\frac{5}{7},\frac{15}{28}\right],\\
\mathcal{U}_{1,6}&=\left[0,\frac{15}{2}\right]\oplus\left[\frac{1}{7},\frac{5}{14}\right]\oplus\left[\frac{5}{7},\frac{39}{14}\right].
\end{aligned}
$$
\end{example}

For $k\geq 2$, $\mathcal{U}_{k}=\mathcal{U}_{k,0}$ is defined by
the commutant VOA of $\mathcal{L}(k+5,0)$ in GKO-construction,
and its modules $\mathcal{U}_{k,i}$ are defined by the summand paired with $\mathcal{L}(k+5,i)$.
%the summand before $\mathcal{L}(k+5,0)$ in GKO-construction,
That is, we write GKO-constructions as
$$\begin{aligned}
&\left(\mathcal{L}\left(3,0\right)\otimes\mathcal{L}\left(1,0\right)\otimes\mathcal{L}\left(1,0\right)\otimes\mathcal{L}\oplus\mathcal{L}\left(3,3\right)\otimes\mathcal{L}\left(1,1\right)\otimes\mathcal{L}\left(1,0\right)\right)\otimes\mathcal{L}(1,0)^{\otimes k}\\
\cong&\bigoplus_{i}\mathcal{U}_{k,i}\otimes \mathcal{L}(k+5,i).
\end{aligned}
$$
For simplicity, we define 
$\left[i-1,h\right]_k:=
\mathcal{U}_{k,i-1}\otimes L\left(c_{k+7},h\right)$ for odd $i$.

Specifically, we have
$$\mathcal{U}_{k,0}=\bigoplus_{i} U_{k-1}^{i}$$
where
$$U_k^{i}=\left[i-1,h_{(1,i)}^{(k+7)}\right]_k=
\mathcal{U}_{k,i-1}\otimes L\left(c_{k+7},h_{(1,i)}^{(k+7)}\right).$$
% $i-1$ travels through all possible $n$
% such that
% $L\left(c_{k+7},h_{(1,i)}^{(k+7)}\right)$ are exactly
% the unitary series appear before 
% $\mathcal{L}(k+5,0)$ in GKO-construction,
% of $\mathcal{L}(1,0)\otimes\mathcal{L}(m,n)$ with $m=k+4$
% and $M_{k}^{(i-1)/2}$ are irreducible modules of $\mathcal{U}_{k-1}$ appear
% in the commutant structure.

% The irreducible modules of $\mathcal{U}_k$ are denoted by $M_{k}^j$.



\begin{example}\label{U2}
We give a precise construction of $\mathcal{U}_2$.
%Collect all $L(c,h)$ appear with
Taking the commutant of $\mathcal{L}(7,0)$,
% in GKO-construction of $\mathcal{L}(1,0)\otimes\mathcal{L}(6,n)$,
we have
%and appearing $L(c,h)$ are 
$L(\frac{11}{12},0)$, $L(\frac{11}{12},\frac{5}{4})$, $L(\frac{11}{12},\frac{19}{4})$, and $L(\frac{11}{12},\frac{21}{2})$ 
with $(m,n)=(6,0),(6,2),(6,4)$, and $(6,6)$.
Similarly as Lemma (\ref{realization of VOA}), we have
\begin{equation*}
\begin{aligned}
&\mathcal{L}\left(3,0\right)\otimes\mathcal{L}\left(1,0\right)\otimes\mathcal{L}\left(1,0\right)\otimes\mathcal{L}\left(1,0\right)\otimes\mathcal{L}\left(1,0\right)\\
&\oplus\mathcal{L}\left(3,3\right)\otimes\mathcal{L}\left(1,1\right)\otimes\mathcal{L}\left(1,0\right)\otimes\mathcal{L}\left(1,0\right)\otimes\mathcal{L}\left(1,0\right)\\
% %\cong&\left\{\left[0,0\right]\oplus\left[1,\frac{5}{4}\right]\oplus\left[2,\frac{19}{4}\right]\oplus\left[3,\frac{21}{2}\right]\right\}\otimes\mathcal{L}\left(7,0\right)\\
% %&\oplus M^0\otimes\left\{L(\frac{11}{12},\frac{7}{9})\otimes\mathcal{L}(7,2)\oplus L(\frac{11}{12},\frac{10}{3})\otimes\mathcal{L}(7,4)\oplus L(\frac{11}{12},\frac{23}{3})\otimes\mathcal{L}(7,6)\right\}\\
% %&\oplus M^1\otimes\left\{L(\frac{11}{12},\frac{1}{36})\otimes\mathcal{L}(7,2)\oplus L(\frac{11}{12},\frac{7}{12})\otimes\mathcal{L}(7,4)\oplus L(\frac{11}{12},\frac{35}{12})\otimes\mathcal{L}(7,6)\right\}\\
% %&\oplus M^2\otimes\left\{L(\frac{11}{12},\frac{55}{36})\otimes\mathcal{L}(7,2)\oplus L(\frac{11}{12},\frac{1}{12})\otimes\mathcal{L}(7,4)\oplus L(\frac{11}{12},\frac{5}{12})\otimes\mathcal{L}(7,6)\right\}\\
% %&\oplus M^3\otimes\left\{L(\frac{11}{12},\frac{95}{18})\otimes\mathcal{L}(7,2)\oplus L(\frac{11}{12},\frac{11}{6})\otimes\mathcal{L}(7,4)\oplus L(\frac{11}{12},\frac{1}{6})\otimes\mathcal{L}(7,6)\right\}\\
\cong&\left\{\left[0,0\right]_1\oplus\left[2,\frac{5}{4}\right]_1\oplus\left[4,\frac{19}{4}\right]_1\oplus\left[6,\frac{21}{2}\right]_1\right\}\otimes\mathcal{L}\left(7,0\right)\\
&\oplus\left\{\left[0,\frac{7}{9}\right]_1\oplus\left[2,\frac{1}{36}\right]_1\oplus\left[4,\frac{55}{36}\right]_1\oplus\left[6,\frac{95}{18}\right]_1\right\}\otimes\mathcal{L}\left(7,2\right)\\
&\oplus\left\{\left[0,\frac{10}{3}\right]_1\oplus\left[2,\frac{7}{12}\right]_1\oplus\left[4,\frac{1}{12}\right]_1\oplus\left[6,\frac{11}{6}\right]_1\right\}\otimes\mathcal{L}\left(7,4\right)\\
&\oplus\left\{\left[0,\frac{23}{3}\right]_1\oplus\left[2,\frac{35}{12}\right]_1\oplus\left[4,\frac{5}{12}\right]_1\oplus\left[6,\frac{1}{6}\right]_1\right\}\otimes\mathcal{L}\left(7,6\right).
\end{aligned}
\end{equation*}

% Let $M^i$ be irreducible $\mathcal{U}_1=\mathcal{U}_{\operatorname{6A}}$-modules with $M^0\cong\mathcal{U}_{6A}$ in 
% (\ref{Real modules}),
% we have
So
$$U_1^{i}=\left[i-1,h_{(1,i)}^{(8)}\right]_1=\mathcal{U}_{1,i-1}\otimes L\left(\frac{11}{12},h_{(1,i)}^{(8)}\right)$$
and we have
$$
\begin{aligned}
\mathcal{U}_{2}=&\mathcal{U}_{2,0}=U_1^1\oplus U_1^3\oplus U_1^5\oplus U_1^7
\\= &\left[0,0\right]_1
\oplus\left[2,\frac{5}{4}\right]_1\oplus\left[4,\frac{19}{4}\right]_1\oplus\left[6,\frac{21}{2}\right]_1,\\
\mathcal{U}_{2,2}=&\left[0,\frac{7}{9}\right]_1\oplus\left[2,\frac{1}{36}\right]_1\oplus\left[4,\frac{55}{36}\right]_1\oplus\left[6,\frac{95}{18}\right]_1,\\
\mathcal{U}_{2,4}=&\left[0,\frac{10}{3}\right]_1\oplus\left[2,\frac{7}{12}\right]_1\oplus\left[4,\frac{1}{12}\right]_1\oplus\left[6,\frac{11}{6}\right]_1,\\
\mathcal{U}_{2,6}=&\left[0,\frac{23}{3}\right]_1\oplus\left[2,\frac{35}{12}\right]_1\oplus\left[4,\frac{5}{12}\right]_1\oplus\left[6,\frac{1}{6}\right]_1.
\end{aligned}
$$
\end{example}

% A series $\mathcal{U}_k$ can be constructed by tensoring $\mathcal{L}\left(1,0\right)$ on the left side of $\cong$.

It is known by Theorem (\ref{uniqueness of 6A-algebra}) that $\mathcal{U}_{1}$ has a unique VOA structure and it satisfies conditions (1) and (2).
%, denoted by $Y$.
For $k\geq 2$, assume $\mathcal{U}_{k-1}$ have a unique VOA structure
and it satisfies conditions (1) and (2).
We will prove that there is a unique VOA structure on $\mathcal{U}_k$ satisfying these conditions.

\begin{remark}\label{rmk}
We simplify $\mathcal{U}_k$, $\mathcal{U}_{k-1,i}$, $L\left(c_{k+7},h_{(1,i)}^{(k+7)}\right)$ and $U_{k-1}^i$ 
by $\mathcal{U}$, $P^i$, $Q^i$ and $U^i$ respectively from now on,
unless otherwise specified.
\end{remark}

If there are no possible ambigious,
we will write $u^a$, or simply $a$, for elements in a module $U^a$.
% for each $a$.
% The corners of $\mathcal{Y}_{b,c}^a$ will be omitted if they are clear in the context.
We use $N,B,\mathcal{Y}$ (and $\tilde{N},\tilde{B},\tilde{\mathcal{Y}}$)
for fusion rules, braiding matrix, and intertwining operator basis
of $P^i$ (and $Q^i$).

%\begin{theorem}
%    The set $\{M^i| i\in I\}$ is closed under the tensor product of category $\mathcal{C}_{V}$. 
%\end{theorem}
\begin{theorem}\label{same fusion rule}
$N_{b,c}^a=\tilde{N}_{b,c}^a$.
\end{theorem}
\begin{proof}
Recall the notations in Theorem (\ref{c1}). Put $M^i=P^{i/2}$, $N^{i'}=Q^{2i-1}$, 
$V_1=\bigoplus_i M^i$, $V_2=\bigoplus_i N^i$,
$V=\mathcal{U}=\bigoplus_i P^i\otimes Q^i$.
Since $Q^{i'}=Q^i$,
% by direct computation of conformal weights,
we have $\tau=\operatorname{Id}$ and hence $N_{b,c}^a=\tilde{N}_{b,c}^a$.
\end{proof}

From this theorem, $P^i$ and $Q^i$ share same fusion rules.
Fix basis $\mathcal{Y}_{b,c}^a$ and $\tilde{\mathcal{Y}}_{b,c}^a$, 
then $\mathcal{I}_{b,c}^a=\mathcal{Y}_{b,c}^a\otimes \tilde{\mathcal{Y}}_{b,c}^a$ is a basis of
$I_{U^1}\ito{U^b}{U^c}{U^a}$.
% $I_{U^1}\begin{pmatrix}
% U^a\\
% U^b&U^c
% \end{pmatrix}.$
So there exist constants $\lambda_{b,c}^a$ such that
$$Y(b\otimes\tilde{b},z)=\sum_{c,a}\lambda_{b,c}^a\mathcal{I}_{b,c}^a(b\otimes\tilde{b},z)=\sum_{c,a}\lambda_{b,c}^a\mathcal{Y}_{b,c}^a(b,z)\otimes\tilde{\mathcal{Y}}_{b,c}^a(\tilde{b},z).$$
We prove $\lambda_{b,c}^a\neq 0$ for all $a,b,c$ compatible with fusion rules in the next section.

\section{$\lambda_{i,j}^k\neq 0$}
When two constant $\lambda$ and $\mu$ are either both non-zero or both zero, 
we will write the relation as $\lambda\sim \mu$.
% then $\lambda_{i,j}^k=0$ if $N_{i,j}^k=0$.
The goal of this section is proving $\lambda_{i,j}^k\sim 1$ for all $i,j,k$ compatible with fusion rules, that is, $\lambda_{i,j}^k\neq 0$ if $N_{i,j}^k=1$.

We denote tuple $(i,j,k,l)$ for 
$$\left\langle {u^i}',Y(u^j,z_1)Y(u^k,z_2)u^l\right\rangle,$$
%and commutativity implies that there is a nonnegative integer $p$
%such that
%$$\left\langle {u^i}',(z_1-z_2)^p Y(u^j,z_1)Y(u^k,z_2)u^l\right\rangle=
%\left\langle {u^i}',(z_1-z_2)^p Y(u^k,z_2)Y(u^j,z_1)u^l\right\rangle,$$
and we write Proposition (\ref{extension property}) as
$(i,j,k,l)\sim(i,k,j,l)$.

Note that the two $\sim$ have different meanings, 
and we can infer the correct one by context.

\subsection{Fundamental Results}
We denote $e_{i,j}$ for $\mathcal{Y}_{c,i}^a(\cdot,z_2)\mathcal{Y}_{b,d}^{i}(\cdot,z_1)\otimes \tilde{\mathcal{Y}}_{c,j}^a(\cdot,z_2)\tilde{\mathcal{Y}}_{b,d}^{j}(\cdot,z_1)$.
Decompose $Y$ to the basis $\{e_{i,j}:i,j=1,\ldots,N\}$,
by direct computation we have
$$
\begin{aligned}
&E\left\langle {(a\otimes \tilde{a})}',
Y(b\otimes \tilde{b},z_1)
Y(c\otimes \tilde{c},z_2)
d\otimes \tilde{d}\right\rangle\\
=&E\left\langle {(a\otimes \tilde{a})}',
\sum_{\beta}\lambda_{b,\beta}^a\lambda_{c,d}^{\beta}
\mathcal{Y}_{b,\beta}^a
\otimes\tilde{\mathcal{Y}}_{b,\beta}^a(b\otimes \tilde{b},z_1)
\mathcal{Y}_{c,d}^{\beta}
\otimes\tilde{\mathcal{Y}}_{c,d}^{\beta}(c\otimes \tilde{c},z_2)
d\otimes \tilde{d}\right\rangle\\
=&E\left\langle {(a\otimes \tilde{a})}',
\sum_{\beta}\lambda_{b,\beta}^a\lambda_{c,d}^{\beta}
\mathcal{Y}_{b,\beta}^a (b,z_1)
\mathcal{Y}_{c,d}^{\beta}(c,z_2)d
\otimes
\tilde{\mathcal{Y}}_{b,\beta}^{a}(\tilde{b},z_1)
\tilde{\mathcal{Y}}_{c,d}^{\beta}(\tilde{c},z_2)
\tilde{d}\right\rangle
\\
=&E\left\langle {(a\otimes \tilde{a})}',
\sum_{\beta}\lambda_{b,\beta}^a\lambda_{c,d}^{\beta}
\sum_{\gamma_1}B_{\beta,\gamma_1}
\mathcal{Y}_{c,\gamma_1}^a(c,z_2)\mathcal{Y}_{b,d}^{\gamma_1}(b,z_1)d
\otimes\sum_{\gamma_2}\tilde{B}_{\beta,\gamma_2}
\tilde{\mathcal{Y}}_{c,\gamma_2}^a(\tilde{c},z_2)\tilde{\mathcal{Y}}_{b,d}^{\gamma_2}(b,z_1)\tilde{d}
\right\rangle
\\
=&E\left\langle {(a\otimes \tilde{a})}',
\sum_{\beta,\gamma_1,\gamma_2}\lambda_{b,\beta}^a\lambda_{c,d}^{\beta}
B_{\beta,\gamma_1}\tilde{B}_{\beta,\gamma_2}
e_{\gamma_1,\gamma_2}d\otimes \tilde{d}\right\rangle
\end{aligned}
$$
and
$$
\begin{aligned}
&E\left\langle {(a\otimes \tilde{a})}',
Y(c\otimes \tilde{c},z_2)Y(b\otimes \tilde{b},z_1)d\otimes \tilde{d}\right\rangle\\
=&E\left\langle {(a\otimes \tilde{a})}',
\sum_{\mu}\lambda_{c,\mu}^a\lambda_{b,d}^{\mu} 
\mathcal{Y}_{c,\mu}^a(c,z_2)\mathcal{Y}_{b,d}^{\mu}(b,z_1)d
\otimes 
\tilde{\mathcal{Y}}_{c,\mu}^a(\tilde{c},z_2)\tilde{\mathcal{Y}}_{b,d}^{\mu}(\tilde{b},z_1)\tilde{d}\right\rangle\\
=&E\left\langle {(a\otimes \tilde{a})}',
\sum_{\mu}\lambda_{c,\mu}^a\lambda_{b,d}^{\mu} 
e_{\mu,\mu}d\otimes\tilde{d}\right\rangle,
\end{aligned}
$$
all possible $\beta,\gamma_1,\gamma_2,\mu$ compatible with fusion rules.

\begin{remark}
Note that elements in braiding matrix is denoted by $\left(B_{a_{4,}a_{1}}^{a_{3},a_{2}}\right)_{\mu,\gamma}^{i,j;k,l}$ in (\ref{braiding matrix property}).
In this paper, $i,j,k,l=1$ by fusion rules, and $a_1,a_2,a_3,a_4$ can always be inferred from the context.
So we can simplify elements in braiding matrix as $B_{\mu,\gamma}$.

Rather than focusing on elements, 
we are going to write the entire braiding matrices.
We need to determine elements by its position in the matrix, 
and we also need to determine the elements by corresponding modules.

To clarify notations, 
we use $B_{i,j}$ to determine the $(i,j)$-element in the matrix,
and $B_{U^i,U^j}$ or $B_{\mu,\gamma}$ the element whose position in the braiding matrix
is determined by (\ref{braiding matrix property}).
\end{remark}

% This euqation can be written as the following matrix form by the following lemma.

\begin{lemma}
% 	Setting $B=\left(B_{i,j}\right)$ and $\tilde{B}=\left(\tilde{B}_{i,j}\right)$,
% we write the equation euqation above as a matrix by iterating all possible pair $(\gamma_1,\gamma_2)$,
% and LHS its non-zero only if $\gamma_1=\gamma_2$, denoted by $\eta$.
% Collecting all possible $\mu$ as $\mu_1,\ldots,\mu_N$ and all $\eta$ as $\eta_1,\ldots,\eta_N$,
% we get
% $$
% \begin{pmatrix}
%     z_{11},&\ldots,&z_{1n}\\
%     \vdots,&&\vdots\\
%     z_{n1},&\ldots,&z_{nn}
% \end{pmatrix}
% \operatorname{diag}(y_1,\ldots,y_n)
% $$
The relation $(a,b,c,d)\sim(a,c,b,d)$
can be written as
\begin{equation}\label{braid matrix and lambda}
B^{\operatorname{T}}\operatorname{diag}\left(\lambda_{b,\beta_1}^a\lambda_{c,d}^{\beta_1},\ldots,\lambda_{b,\beta_N}^a\lambda_{c,d}^{\beta_N}\right)\tilde{B}
=\operatorname{diag}\left(\lambda_{c,\mu_1}^a\lambda_{b,d}^{\mu_1},\ldots,\lambda_{c,\mu_N}^a\lambda_{b,d}^{\mu_N}\right). 
\end{equation}
where $B=\left(B_{i,j}\right)$,
$\tilde{B}=\left(\tilde{B}_{i,j}\right)$,
and $\mu_i, \beta_j \left(i,j=1,\ldots,N\right)$ are all modules compatible with fusion rules.
\end{lemma}
\begin{proof}
Let $C=\operatorname{diag}\left(\lambda_{b,\beta_1}^a\lambda_{c,d}^{\beta_1},\ldots,\lambda_{b,\beta_N}^a\lambda_{c,d}^{\beta_N}\right)$,
$D=\operatorname{diag}\left(\lambda_{c,\mu_1}^a\lambda_{b,d}^{\mu_1},\ldots,\lambda_{c,\mu_N}^a\lambda_{b,d}^{\mu_N}\right)$,
and $A=B^{\operatorname{T}}C\tilde{B}$.
Now we check the $A_{i,j}=D_{i,j}$.
By direct computation we know
$$A_{i,j}=\sum_{n,m=1}^N B_{n,i}C_{n,m}\tilde{B}_{m,j}=\sum_{k=1}^N B_{k,i}C_{k,k}\tilde{B}_{k,j}.$$ 

Consider $(a,b,c,d)\sim(a,c,b,d)$.
Use the linear independence of four point functions,
we get 
$$
\sum_{\beta,\gamma_1,\gamma_2}\lambda_{b,\beta}^a\lambda_{c,d}^{\beta}
B_{\beta,\gamma_1}\tilde{B}_{\beta,\gamma_2}
e_{\gamma_1,\gamma_2}
=\sum_{\mu}\lambda_{c,\mu}^a\lambda_{b,d}^{\mu} 
e_{\mu,\mu}.
$$
LHS of it are non-zero only if $\gamma_1=\gamma_2=\mu$.
%denoted by $\eta_j,j=1,\ldots,N$.
Since $e_{i,j}$ are basis,
we have 
\begin{equation}\label{DJY equations}
	\sum_{\beta}\lambda_{b,\beta}^a\lambda_{c,d}^{\beta}B_{\beta,\mu_i}\tilde{B}_{\beta,\mu_i}
=\lambda_{c,\mu_i}^a\lambda_{b,d}^{\mu_i}
\end{equation}
for all possible $\mu_i$ compatible with fusion rules.

So we get $A_{i,j}=0=D_{i,j}$ if $i\neq j$.
%And there exsit permutaions $\sigma(\eta_i)=\mu_i$, 
And we have
$$D_{i,i}=\lambda_{c,\mu_i}^a\lambda_{b,d}^{\mu_i}
=\sum_{\beta}\lambda_{b,\beta}^a\lambda_{c,d}^{\beta}B_{\beta,\mu_i}\tilde{B}_{\beta,\mu_i}
=\sum_{k=1}^N C_{k,k}B_{k,i}\tilde{B}_{k,i}=A_{i,i},$$
thus completing the proof.
\end{proof}
%$$\lambda_{b,\beta}^a\lambda_{c,d}^{\beta}
%B_{\beta,\mu_p}\tilde{B}_{\beta,\mu_p}=\lambda_{c,\mu_p}^a\lambda_{b,d}^{\mu_p} $$
%for all possible $\mu_p$ compatible with fusion rules.


%But the linear independence in assert that 

% Note that the linear independence of four point functions makes elements non-zero only if $\gamma_1=\gamma_2$
	
% LHS of (\ref{linear independence}) can be written as a matrix by iterating all possible pair $(\gamma_1,\gamma_2)$,
% and they are non-zero only if $\gamma_1=\gamma_2$, denoted by $\eta_p,p=1,\ldots,M$.


\begin{remark}
	Instead of (\ref{DJY equations}),
	\cite{DJY} use system of equations in Theorem 3.9.
\end{remark}

%Take trace on both side of (\ref{braid matrix and lambda}),
%Take trace of both sides of (\ref{braid matrix and lambda}),
%
From the previous lemma
we get
\begin{equation}\label{braid matrix and lambda sum}
\left(\lambda_{c,\mu_1}^a\lambda_{b,d}^{\mu_1},\ldots,\lambda_{c,\mu_N}^a\lambda_{b,d}^{\mu_N}\right)
=\left(\sum_{j=1}^N\lambda_{b,\beta_j}^a\lambda_{c,d}^{\beta_j}B_{\beta_j,\mu_1}\tilde{B}_{\beta_j,\mu_1},\ldots,\sum_{j=1}^N\lambda_{b,\beta_j}^a\lambda_{c,d}^{\beta_j}B_{\beta_j,\mu_N}\tilde{B}_{\beta_j,\mu_N}\right)
\end{equation}

% \begin{equation}\label{braid matrix and lambda sum}
% \sum_{\alpha} \prod_{\beta} \lambda_{(\alpha,\beta)}B_{(\alpha,\beta)}\tilde{B}_{(\alpha,\beta)}
% =\sum_{\alpha'} \prod_{\beta'} \lambda_{(\alpha',\beta')},
% \end{equation}
% where $(\alpha,\beta)$ and $(\alpha',\beta')$ give indices for every possible braiding matrix $B,\tilde{B}$ and $\lambda_{i,j}^k$ compatible with fusion rules.


%We write this relation of $\lambda_{i,j}^k$ as 
Since we don't need the order of $\mu_i$ and $\beta_j$,
and the values of $B_{i,j}$ and $\tilde{B}_{i,j}$ are irrelevant in this section,
we write the equation (\ref{braid matrix and lambda sum}) which comes fom $(a,b,c,d)\sim(a,c,b,d)$ 
as a relation between two sets of numbers:
\begin{equation}\label{lambda equation}
\left\{\lambda_{c,\mu_1}^a\lambda_{b,d}^{\mu_1},\ldots,\lambda_{c,\mu_N}^a\lambda_{b,d}^{\mu_N}\right\}
\sim\left\{\lambda_{b,\beta_1}^a\lambda_{c,d}^{\beta_1},\ldots,\lambda_{b,\beta_N}^a\lambda_{c,d}^{\beta_N}\right\}.
\end{equation}

\begin{theorem}\label{samerank}
The none-zero elements appear same times in each side of (\ref{lambda equation}).
\end{theorem}
\begin{proof}
Noting that $B,\tilde{B}$ are invertible and two sides of (\ref{braid matrix and lambda}) share the same rank, the proof is immediate.
\end{proof}

\begin{remark}
Note that if there is only one element in (\ref{lambda equation}), this relation will be compatible with the $\sim$ between numbers.
In this case we will omit the braces.
\end{remark}

In order to get the uniqueness of $\mathcal{U}_k$,
we need to know some elements in $B$ and $\tilde{B}$.
We will show in the next section that $B$ can by computed by $\tilde{B}$,
and non-zero elements $\tilde{B}_{i,j}$ can be checked cace by case for any fixed $\mathcal{U}_k$.

Now we are ready to start proving $\lambda_{i,j}^k\sim 1$ for all $i,j,k$ compatible with fusion rules.

\begin{lemma}
$\lambda_{1,k}^{k}\neq 0$,
$\lambda_{k,1}^{k}\neq 0$ 
,and $\lambda_{k,k}^{1}\neq 0$.
\end{lemma}
\begin{proof}
The first claim is a direct consequence from the definition that $U^k$ is an irreducible $U^1$-module.

For any $u^{k}\in U^k$, using skew symmetry of $Y\left(\cdot,z\right)$,
we have
\[
Y(u^{k},z)u^{1}=e^{zL\left(-1\right)}Y\left(u^{1},-z\right)u^{k}=\lambda_{1,k}^{k}\cdot e^{zL\left(-1\right)}\mathcal{I}_{1,k}^{k}\left(u^{1},-z\right)u^{k}=\lambda_{k,1}^{k}\cdot\mathcal{I}_{k,1}^{k}\left(u^{k},z\right)u^{1}.
\]
So $\lambda_{k,1}^{k}\not=0$.

For $u^{k},v^{k}\in U^k$, we have
\[
\left\langle Y\left(u^{k},z)v^{k}\right),u^{1}\right\rangle =\left\langle v^{k},Y\left(e^{zL\left(-1\right)}\left(-z^{-2}\right)^{L\left(0\right)}u^{k},z^{-1}\right)u^{1}\right\rangle .
\]
That is,
\[
\left\langle \lambda_{k,k}^{1}\cdot\mathcal{I}_{k,k}^{1}\left(u^{k},z\right)v^{k},u^{1}\right\rangle =\left\langle v^{k},\lambda_{k,1}^{k}\cdot\mathcal{I}_{k,1}^{k}\left(e^{zL\left(-1\right)}\left(-z^{-2}\right)^{L\left(0\right)}u^{k},z^{-1}\right)u^{1}\right\rangle .
\]
Applying previous claim, $\lambda_{k,1}^{k}\not=0$, and hence $\lambda_{k,k}^{1}\not=0$.
\end{proof}

\begin{theorem}\label{alternation}
$\lambda_{i,j}^k\sim\lambda_{j,i}^k\sim\lambda_{j,k}^i$ for all $i,j,k$ compatible with fusion rules.
\end{theorem}
\begin{proof}
Consider $(1,i,j,k)\sim (1,j,i,k)$.
Note that there will be only one monomial in each side of the equation
since $N_{a,b}^1\neq 0$ if and only if $a=b$.
We have
%$$
%\lambda_{j,k}^i\lambda_{i,i}^1\left\langle u^1,\mathcal{I}_{i,i}^1(u^i,z_1)\mathcal{I}_{j,k}^i(u^j,z_1)u^k\right\rangle
%=\lambda_{i,k}^j\lambda_{j,j}^1\left\langle u^1,\mathcal{I}_{j,j}^1(u^j,z_1)\mathcal{I}_{i,k}^j(u^i,z_1)u^k\right\rangle.
%$$
%So 
$\lambda_{j,k}^i\lambda_{i,i}^1=\lambda_{i,k}^j\lambda_{j,j}^1.$
So $\lambda_{j,k}^i\sim \lambda_{i,k}^j$ by previous claim.
\end{proof}

Transpositions of indices of $\lambda_{a,b}^c$ will be omitted from now on.

%Fusion rules show that the longest direct summand appears at the $2K+1$st position when $m=4K+1$ and at the $2K+3$rd position when $m=4K+3$.
%When $m=4K$, the longest direct summand appears at the $2K+1$st position. 
%When $m=4K+2$, the longest direct summand appear at both the $2K+1$st and $2K+3$rd position.

% \begin{lemma}
%     $\lambda_{3,3}^5\sim 1$.
% \end{lemma}
% \begin{proof}
%     Assume $\lambda_{3,3}^5=0$, 
%     then $Y(u^{3},z)U^3\subseteq U^1+U^3$ which makes $U^1+U^3$ a subVOA of $\mathcal{U}$.
%     Since mirror extension theory asserts that the category of irreducible modules $\mathcal{C}_{V}$ for $V=U^1+U^3$ is closed under tensor product,
%     it contradicts the fusion rule of $U^3\boxtimes U^3$.
% \end{proof}


\begin{lemma}\label{3,i,i-2}
$\lambda_{i,i-2}^3\sim 1$ for all $i$ compatible with fusion rules.
%$3\leq i\leq 2K+3$.
\end{lemma}
\begin{proof}
We already know $\lambda_{3,1}^3\sim 1$. 
Assume $\lambda_{3,i-2}^i\sim 1$ for $i=3, 5,7,\dots,2t+1$,
but $\lambda_{3,2t+1}^{2t+3}=0$.

If we can prove that $\lambda_{p,q}^{r}=0$ for all 
$1\leq \{p,q\}\leq 2t+1,r\geq 2t+3$ when $N_{p,q}^r=1$, 
then $V=\sum_{n=1}^{t}U^{2n+1}$ will be a subVOA of $\mathcal{U}$.
% Since mirror extension theory asserts that the category of irreducible modules $\mathcal{C}_{V}$ is closed under tensor product,
It has a contradiction between Theorem (\ref{c1})(ii) and the fusion rule of $U^{2n+1}\boxtimes U^{2n+1}$ for some $n$ with $U^{2t+3}$ being a summand, 
which makes our assumption impossible.

To get the contradiction we need, 
we will firstly prove that $\lambda_{p,q}^{r}=0$ for all $r=2t+3$ and $p,q$ such that $p+q=2t+4$.
Then we will prove that $\lambda_{p,q}^r=0$ for all $r\geq 2t+3,p\leq 2t-1,q\leq 2t+1$.
Thus we can get the contradiction.
% by induction on $p$ firstly.

For $r=2t+3$ and $p+q=2t+4$,
we consider $(3,p,q-2,r)\sim (3,q-2,p,r)$.
And we get
%\begin{equation}\label{(3,p,q-2,2t+3)}
$$
\left\{
\lambda_{3,p}^{p-2}\lambda_{r,q-2}^{p-2},\quad
\lambda_{3,p}^{p}\lambda_{r,q-2}^{p},\quad
\lambda_{3,p}^{p+2}\lambda_{r,q-2}^{p+2}
\right\}
\sim
\left\{
\lambda_{3,q-2}^{q-4}\lambda_{r,p}^{q-4},\quad
\lambda_{3,q-2}^{q-2}\lambda_{r,p}^{q-2},\quad
\lambda_{3,q-2}^{q}\lambda_{r,p}^{q}\right\}.
$$
%\end{equation}
Note that by fusion rules we have
% $N^{i}_{a,b}=0$ if $|a-b|> 2i-2$,
% we get 
$N_{r,q-2}^{p-2}=N_{r,q-2}^{p}=N_{r,p}^{q-4}=N_{r,p}^{q-2}=0$,
and that our assumption asserts $\lambda_{3,i-2}^{i}\sim 1$ for odd $i$ such that $3\leq i \leq 2t+1$.
As a result, the equation becomes
%$$\lambda_{3,p}^{p+2}\lambda_{2t+3,2t+2-p}^{p+2}
%\sim\lambda_{3,2t+2-p}^{2t+4-p}\lambda_{2t+3,p}^{2t+4-p}.$$
$\lambda_{r,q-2}^{p+2}
\sim\lambda_{r,p}^{q}$ for all odd $p\leq 2t+1$.
So 
$$\lambda_{r,p}^{q}
\sim\lambda_{r,p-2}^{q+2}\sim\cdots\sim \lambda_{r,3}^{2t+1}=0,$$
and the first step is finished.

%Let $p=3$, and we have

%$$\lambda_{2t+3,2t-1}^{1}+\lambda_{3,3}^3\lambda_{3,2t+3}^{2t-1}+\lambda_{5,2t+3}^{2t-1}
%\sim\lambda_{3,2t+3}^{2t-3}+\lambda_{3,2t-1}^{2t-1}\lambda_{2t+3,3}^{2t-1}+\lambda_{2t+3,3}^{2t+1}.$$
%Note that $N_{i,j}^3=0$ if $|i-j|>2$. This leaves $\lambda_{5,2t+3}^{2t-1}=0$ by Theorem (\ref{samerank}).

%Assume we already get what we want when $3\leq p\leq m$. 
%Now take $p=m+2$.
%Since $N_{2t+3,p-2}^{2t+2-p}=N_{2t+3,p}^{2t-p}=N_{2t+3,p}^{2t+2-p}=0$,
%we have $\lambda_{2t+3,m+2}^{q}=\lambda_{2t+3,p}^{2t+4-p}\sim\lambda_{2t+3,2t+2-p}^{p+2}=0$

Now we consider the case when $p\leq 2t-1,q\leq 2t+1,r\geq 2t+3$.
%It is shown by fusion rule that $N_{p,q}^r=0$ for small $p+q$.
Pick $p,q$ such that $N_{p',q'}^r=0$ for every possible $p'\leq p, q'\leq q$ and the sum $p+q$ is the biggest.
%the following statements claim that $\lambda_{p+2,q}^r=0$.
%Thus $\lambda_{p,q}^r=0$ for all $p,q\leq 2t+1$.

Consider $(3,p,q,r)\sim(3,q,p,r)$,
we have
$$\left\{
	\lambda_{3,p}^{p-2}\lambda_{q,r}^{p-2},\quad
\lambda_{3,p}^{p}\lambda_{q,r}^{p},\quad
\lambda_{3,p}^{p+2}\lambda_{q,r}^{p+2}\right\}
\sim
\left\{
\lambda_{3,q}^{q-2}\lambda_{p,r}^{q-2},\quad
\lambda_{3,q}^{q}\lambda_{p,r}^{q},\quad
\lambda_{3,q}^{q+2}\lambda_{p,r}^{q+2}\right\}.$$
It also holds that $\lambda_{p-2,q}^{r}=\lambda_{p,q}^{r}=\lambda_{p,q-2}^{r}=0$ by assumption.
Since $p\leq 2t-1$, $\lambda_{3,p}^{p+2}\sim 1$,
we get $\lambda_{q,r}^{p+2}\sim\lambda_{3,q}^{q+2}\lambda_{p,r}^{q+2}$.

If $q=2t+1$, then by assumption $\lambda_{q,r}^{p+2}\sim\lambda_{q,3}^{q+2}=0$.
So we proved $\lambda_{p,2t+1}^r=0$ for any $1\leq p\leq 2t+1$.
If $q=2t-1$, then $\lambda_{q,r}^{p+2}\sim\lambda_{p,r}^{q+2}=\lambda_{p,r}^{2t+1}=0$.
We get $\lambda_{p,q}^r=0$ for any $p,q\leq 2t+1$ by induction on $q$.
% $(3,3,2t-1,2t+3)\sim(3,2t-1,2t+3)$ implies $\lambda_{5,2t-1}^{2t+3}=0$.
Thus the second step is finished and we get the contradiction we need.
\end{proof}

\begin{lemma}{\label{recursion}}
Assume for any odd $i$, $\lambda_{i,i}^3\sim 1$.
For any odd $p$, if $\lambda_{i,j}^k\sim 1$ holds as long as
$\min\left\{i,j,k\right\}\leq p$
with $i,j,k$ compatible with fusion rules,
then $\lambda_{i,j}^{p+2}\sim 1$.
\end{lemma}
%\left\{1,3,\dots, 2p+1\right\}\neq \emptyset$,

\begin{proof}
$(3,p,i,j)\sim(3,i,p,j)$ implies
$$
\left\{\lambda_{3,p}^{p-2}\lambda_{i,j}^{p-2},\quad
\lambda_{3,p}^{p}\lambda_{i,j}^{p},\quad
\lambda_{3,p}^{p+2}\lambda_{i,j}^{p+2}\right\}
\sim
\left\{
\lambda_{3,i}^{i-2}\lambda_{p,j}^{i-2},\quad
\lambda_{3,i}^{i}\lambda_{p,j}^{i},\quad
\lambda_{3,i}^{i+2}\lambda_{p,j}^{i+2}\right\}.
$$
Using assumptions and the previous lemma, we get $\lambda_{i,j}^{p+2}\sim 1$.
\end{proof}

Now if we can prove $\lambda_{i,i}^3\sim 1$ for all $i$ compatible with fusion rules, 
we will get $\lambda_{i,j}^k\neq 0$ for all $i,j,k$ compatible with fusion rules by induction using the previous lemma.
In order to prove this fact, 
we need to consider two cases separately:
the parity of the biggest index, 
denoted by $m$, 
appears in fusion rules.

\subsection{When $m$ is odd}
The case when $m$ is assumed to be odd is under our first consideration.

% We have $U^3\boxtimes U^3=U^1+U^3+U^5$ by fusion rule.
%It can also be checked directly that $U^m\boxtimes U^m=U^1+U^3$ for $m$ even and $U^{m+1}\boxtimes U^{m+1}=U^{1}$ otherwise.

\begin{lemma}
$\lambda_{m,m}^3\sim 1$.
\end{lemma}
\begin{proof}
It can be checked directly that $U^m\boxtimes U^m=U^1+U^3$.
If $\lambda_{m,m}^3=0$, then $V=U^1+U^m$ is a subVOA of $\mathcal{U}$.
The contradiction rises since $U^m\boxtimes U^m$ was supposed to be closed under tensor products in $\mathcal{C}_V$ by Theorem (\ref{c1})(ii). 
%    The claim holds for parallel lines in previous lemma with proper subsititutions.
\end{proof}

% By considering locality  $(i,j,k,l)\sim (i,k,j,l)$ in each case,
% noting that none-zero monomials of modules' tensor products appear same times in each side of $\sim$,
% direct results for some special cases are shown as follwing:

% % $$(m,m,3,5)\Rightarrow\lambda_{m,m-2}^3\neq 0,\lambda_{m,m-2}^5\neq 0,$$

% % $$(m,m,3,3)\Rightarrow\lambda_{3,3}^3\neq 0,$$

% % $$(m,m,m-2,m-2)\Rightarrow\lambda_{m-2,m-2}^3\neq 0,$$

% % $$(m,m,5,5)\Rightarrow\lambda_{5,5}^3\sim\lambda_{m,m-4}^5,$$

% % $$(m,m,m-2,m-4)\Rightarrow\lambda_{m-2,m-4}^3\sim\lambda_{m,m-4}^5.$$

% $$\begin{aligned}
%     (m,m,3,5)&\Rightarrow\lambda_{m,m-2}^3\sim 1,\lambda_{m,m-2}^5\sim 1,\\
%     (m,m,3,3)&\Rightarrow\lambda_{3,3}^3\sim 1,\\
%     (m,m,m-2,m-2)&\Rightarrow\lambda_{m-2,m-2}^3\sim 1,\\
%     (m,m,5,5)&\Rightarrow\lambda_{5,5}^3\sim\lambda_{m,m-4}^5,\\
%     (m,m,m-2,m-4)&\Rightarrow\lambda_{m-2,m-4}^3\sim\lambda_{m,m-4}^5,\\
%     (m,m,i,i)&\Rightarrow\lambda_{m,m}^3\lambda_{i,i}^3+\lambda_{m,m}^1\lambda_{i,i}^1\sim\left(\lambda_{m,i}^{m-i+1}\right)^2+\left(\lambda_{m,i}^{m-i+3}\right)^2.
% \end{aligned}$$

% \begin{example}
%     $\lambda_{3,5}^7\sim 1.$
% \end{example}

% \begin{proof}
% When $m\geq 9$, $(7,5,3,3)$ and $(9,5,3,3)$ give
% $$
% \begin{cases}
%     \lambda_{7,5}^3\lambda_{3,3}^3+\lambda_{3,3}^5\lambda_{5,5}^7\sim \lambda_{7,3}^5\lambda_{5,5}^3+\lambda_{7,7}^3\lambda_{5,7}^3\\
%     \lambda_{5,5}^9\lambda_{3,3}^5\sim \lambda_{9,3}^7\lambda_{3,5}^7
% \end{cases}$$

% It's suffice to prove by contradiction. Assume $\lambda_{3,5}^7=0$, then $\lambda_{5,5}^9=\lambda_{5,5}^7=0$.
% These conditions make $U^1+U^3+U^5$ a subVOA of $\mathcal{U}$.
% This is impossible by the same statement as in the previous lemma with replacement $V=U^1+U^3+U^5$.
% \end{proof}


% %\subsection{ $\lambda_{i,j}^k\sim 1$ for every $N_{i,j}^k\neq 0$}
% We need $\lambda_{i,j}^k\sim 1$ for every $N_{i,j}^k\neq 0$. 
% Several lemmas are needed for this claim.


%Firstly we discuss under the assumption that
%$\lambda_{i,i-2}^3\sim 1$ hold for every odd $i$
%when $3\leq i\leq 2K+3$.

\begin{lemma}
$\lambda_{a,b}^m\sim 1$ for all $a,b$ compatible with fusion rules.
\end{lemma}
\begin{proof}
By $(m,i,m,i-2)\sim(m,m,i,i-2)$ we have
$$\lambda_{i,m-i+3}^m\lambda_{m,i-2}^{m-i+3}\sim\lambda_{m,m}^3\lambda_{i,i-2}^3\sim 1,$$
then
$\lambda_{m,i}^{m-i+3}\sim \lambda_{m,i-2}^{m-i+3}\sim 1$
and these are all cases of $a,b$ that is compatible with fusion rules.
% $$\lambda_{m,m-i+3}^m\sim\cdots\sim\lambda_{m,m-i+3}^i\sim \lambda_{m,m-i+3}^{i-2}\sim\cdots\sim \lambda_{m,m-i+3}^{1}\sim 1,$$
% thus $\lambda_{a,b}^m\sim 1,\forall a,b$.
\end{proof}
\begin{lemma}
$\lambda_{a,b}^3\sim 1$ for all $a,b$ compatible with fusion rules.
\end{lemma}
\begin{proof}
Consider $(m,m,i,j)$ and we have
$$
\left\{\lambda_{m,m}^1\lambda_{i,j}^1,\quad
\lambda_{m,m}^3\lambda_{i,j}^3\right\}
\sim
\left\{\lambda_{m,i}^{m+1-i}\lambda_{m,j}^{m+1-i},\quad
\lambda_{m,i}^{m+3-i}\lambda_{m,j}^{m+3-i}\right\}.$$
Applying previous lemmas, the proof is direct.
% gives
%$$1+\lambda_{i,j}^3\sim$$
%    Note that elements in $U^m$ appear in $(m,m,i,j)$ only if $j=i-2,i,i+2$,
%while the first two cases have been discussed already,
%in particular, $\lambda_{i,i}^3\sim 1$.
%When $j=i+2$,
%$\lambda_{m,m}^3\lambda_{i,i+2}^3\sim \lambda_{i,m-i+1}^m\lambda_{m,i+2}^{m-i+1}$,
%and we have $\lambda_{i,i+2}^3\sim 1$.
%This shows $\lambda_{a,b}^3\sim 1,\forall a,b$.
\end{proof}

% \begin{lemma}
%     $\lambda_{a,b}^5\sim 1,\forall a,b$.
% \end{lemma}
% \begin{proof}
%     Consider $(3,3,i,j)\sim (3,i,3,j)$, we have
% $$\lambda_{i,j}^1+\lambda_{i,j}^3+\lambda_{i,j}^5\sim \lambda_{i-2,j}^3+\lambda_{i,j}^3+\lambda_{i+2,j}^3\sim1+1+1$$
% (the cases $i=1,m$ are easy to check otherwise),
% thus $\lambda_{i,j}^5\sim 1$.
% \end{proof}


%Now we check the assumption before the lemma 1.7.

%And we get from above lemmas that
\begin{theorem}
If $m$ is odd, then $\lambda_{i,j}^k\neq 0$ for every $N_{i,j}^k =1$.
\end{theorem}
\begin{proof}
%    Since we have proved that $\lambda_{3,3}^3\sim 1$,
%    Use the previous lemma and (\ref{recursion}) with proper alternations and the theorem is archived.
This is a direct consequence of previous lemma and Lemma (\ref{recursion}) with proper
alternations of indices.
\end{proof}

\subsection{When $m$ is even}
% Since $U^{m+1}$ is a simple current module, 
% From $(m-k+2,m+1,i,j)\sim (m-k+2,i,m+1,j)$ we get
% $$\lambda_{m-k+2,m+1}^k\lambda_{i,j}^k\sim\lambda_{m-k+2,i}^{m-j+2}\lambda_{m+1,j}^{m-j+2}.$$
% Thus 
% $\lambda_{i,j}^k\sim\lambda_{i,m-j+2}^{m-k+2}.$
% Note that it also holds after corner alternations. 
% Furthermore, put $i=j$ and we get $\lambda_{i,i}^k\sim\lambda_{m+2-i,m+2-i}^k$. 
Thanks to Lemma (\ref{recursion}), 
we only need
$\lambda_{i,i}^3\sim 1$
whenever $N_{i,i}^3=1$.
% for all $i$ compatible with fusion rules.
% for any odd $i\leq 2K+1$.
%and $\lambda_{i,i+2}^3\sim 1$ for $i\leq 2K+1$.
%When $m$ is even, we have $\mathcal{U}=U^1+U^3+\cdots+U^{m+1}$.
%While the $\lambda_{i,i+2}^3\sim 1$ part is parallel to previous section, only the former condition is needed.
Several lemmas are needed to establish the proof
since $N_{m+1,m+1}^3=0$.

\begin{lemma}\label{5,i,i}
$\lambda_{i,i}^5\sim 1$, $\forall i$ when $N_{i,i}^5=1$.
\end{lemma}
\begin{proof}
Consider $(3,3,i,i)$, $i\geq 3$.
We get 
$$
\left\{1,\quad
\lambda_{3,3}^3\lambda_{i,i}^3,\quad
\lambda_{3,3}^5\lambda_{i,i}^5\right\}
\sim 
\left\{1,\quad
\left(\lambda_{i,i}^3\right)^2,\quad
1\right\}.$$
Thus $\lambda_{i,i}^5\sim 1$.
\end{proof}

\begin{lemma}\label{5,i,i+4}
$\lambda_{i,i+4}^5\sim 1$, $\forall i$ when $N_{i,i+4}^5=1$.
\end{lemma}
\begin{proof}
This comes from $(3,3,i,i+4)$:
$$
\left\{\lambda_{3,3}^1\lambda_{i,i+4}^1,\quad
\lambda_{3,3}^3\lambda_{i,i+4}^3,\quad
\lambda_{3,3}^5\lambda_{i,i+4}^5\right\}
\sim
\left\{
\lambda_{3,i}^{i-2}\lambda_{3,i+4}^{i-2},\quad
\lambda_{3,i}^{i}\lambda_{3,i+4}^{i},\quad
\lambda_{3,i}^{i+2}\lambda_{3,i+4}^{i+2}\right\}.
$$
Since $N_{i,i+4}^1=N_{i,i+4}^3=N_{3,i+4}^{i-2}=N_{3,i+4}^{i}=0$ by fusion rules,
and we proved Lemma (\ref{3,i,i-2}),
the equation becomes $\lambda_{i,i+4}^5\sim 1$.
\end{proof}

% \begin{lemma}
%     $\lambda_{3,3}^3\sim 1$.
% \end{lemma}
% \begin{proof}
%     (?)
% %Assume $\lambda_{3,3}^3=0$.
% %    Recall $(3,3,i,i)$, $i\geq 3$,
% %    and this make $\lambda_{i,i}^3=0, \forall i\geq 3$.
% \end{proof}

\begin{lemma}
If either $\lambda_{3,3}^3=0$ or $\lambda_{5,5}^3=0$,
then $\lambda_{4i+1,4i+1}^3=0$ and $\lambda_{5,i+2}^i=0$, $\forall i$.
\end{lemma}
\begin{proof}
Consider $(3,3,i,i+2)$ and $(3,5,i,i+2)$,
and we get 
\begin{equation}\label{(3,3,i,i+2)}
\begin{aligned}
	\left\{
\lambda_{3,3}^3,\quad\lambda_{i,i+2}^{5}\right\}&\sim
\left\{\lambda_{3,i}^{i},\quad
\lambda_{3,i+2}^{i+2}\right\},\\
\left\{\lambda_{3,5}^5\lambda_{i,i+2}^{5},\quad
\lambda_{i,i+2}^{7}\right\}
&\sim
\left\{\lambda_{3,i}^{i}\lambda_{5,i+2}^{i},\quad
1\right\}.
\end{aligned}
\end{equation}

%The second equation asserts that $\lambda_{3,3}^3$ and $\lambda_{5,5}^3$
%can't be $0$ at same time.


The case $\lambda_{5,5}^3=0$:
This make
$\lambda_{3,i}^{i}\lambda_{5,i+2}^{i}=0$.
% (and $\lambda_{i,i+2}^{7}\sim 1$), $\forall i$.
We can multiply $\lambda_{3,i}^i$ on both side of the first equation.
% Together with %(\ref{(3,3,i,i+2)}),
If $\lambda_{3,3}^3=0$, 
we get $\lambda_{3,i}^{i}=0$ for all $i$,
thus $\lambda_{5,i+2}^i=0$, $\forall i$.
% once there exists $k$ such that $\lambda_{5,k+2}^{k}\sim 1$,
% we have
% $$\lambda_{5,k+2}^{k}\sim 1\Rightarrow
% \lambda_{3,k}^{k}=0\Rightarrow
% \lambda_{3,k+2}^{k+2}\sim 1\Rightarrow
% \lambda_{5,k+2}^{k+4}=0\Rightarrow
% \lambda_{3,k+2}^{k+2}=0,$$
% and the contradiction arises.
% the only way to keep two equations is making $\lambda_{3,i}^{i}=0$ and $\lambda_{5,i+2}^i=0$, $\forall i$.
% and $\lambda_{5,5}^3=0$,
If $\lambda_{3,3}^3\sim 1$, 
we have $\lambda_{4t-1,4t-1}^3\sim 1$ and $\lambda_{4t+1,4t+1}^3=0$
for any $t\geq 1$.
Furthermore, $\lambda_{5,i+2}^i=0$, $\forall i$.

The case $\lambda_{3,3}^3=0$:
Put $i=3$ in the second equation and we get $\lambda_{5,5}^3=0$.
It has been discussed already.
\end{proof}


% And $(3,5,i,i)$ gives
% $$\lambda_{i,i}^{7}\sim
% \lambda_{5,i}^{i-2}+
% \lambda_{5,i}^{i+2}.$$

\begin{lemma}
$\lambda_{3,3}^3\sim\lambda_{5,5}^3\sim 1$.
\end{lemma}
\begin{proof}
Otherwise if we can proof $\lambda_{4p+1,4q+1}^{4r+3}=0$ for all $p,q\geq 1$ and $r\geq 0$,
we will get a subVOA $V$ of $\mathcal{U}$ where $V=\sum_{n=1}^t U^{4n+1}$.
Again there is an $n$ such that the fusion rule of $U^n\boxtimes U^n$ is contradicted to the fact that $\mathcal{C}_V$ is closed under tensor products.

%Assume $\lambda_{5,5}^3=0$.
%According to the previous lemma, 
%it's suffice to prove the whole claim if there are contradictions when $\lambda_{3,3}^3=0$ and $\lambda_{3,3}^3\sim 1$.

It's enough to assume $\lambda_{5,5}^3\sim 0$ and use previous lemma to assert the contradiction.

%  Consider $(5,3,i,i+4)$, we get
% $$\lambda_{i,i+4}^{7}\sim\lambda_{3,i+4}^{i+4}$$

% Consider $(5,3,4p+1,4q+1)$, we have
% $$\lambda_{4p+1,4q+1}^3+\lambda_{4p+1,4q+1}^5+\lambda_{4p+1,4q+1}^7\sim
% \lambda_{5,4p+1}^{4q-1}+\lambda_{5,4p+1}^{4q+1}\lambda_{3,4q+1}^{4q+1}+\lambda_{5,4p+1}^{4q+3}.$$

%When $m=4K+2$,
%contradiction arises since $\lambda_{3,3}^3\sim \lambda_{m-1,m-1}^3$.
%When $m=4K$,
%contradiction arises since $\lambda_{m+1,m+1}^3\sim 1$.

%Let's check the case $m=4K$.

%Assumption makes $\lambda_{i,i}^7\sim \lambda_{5,i}^{i-2}+\lambda_{5,i}^{i+2}$.
%so $\lambda_{i,i}^7=0$, $\forall i\geq 3$.


% Consider $(3,5,i,i+4)\sim (3,i,5,i+4)$, then we have
% $$\lambda_{3,5}^{5}\lambda_{i,5}^{i+4}+
% \lambda_{3,5}^{7}\lambda_{i,i+4}^{7}
% \sim \lambda_{i,i}^{3}\lambda_{i,i+4}^{5}+
% \lambda_{3,i}^{i+2}\lambda_{i+4,i+2}^{5}.$$
% And
% $$\lambda_{i,i+4}^{7}\sim\lambda_{i,i}^{3}\sim\begin{cases}
%     0,&i=4k+1\\
%     1,&i=4k+3.
% \end{cases}$$

Our origion assumption asserts the case $r'=0,p'=1,q'=1$.
We need to show that if $\lambda_{4p+1,4q+1}^{4r+3}=0$ for every $p\leq p',q\leq q',r\leq r'$,
then $\lambda_{4p'+1,4q'+1}^{4(r'+1)+3}=0$ and  $\lambda_{4(p'+1)+1,4q'+1}^{4r'+3}=0$.

Consider $(5,4r+3,4p+1,4q+1)$,
we get
$$\left\{\begin{aligned}
&\lambda_{5,4r+3}^{4r-1}\lambda_{4p+1,4q+1}^{4(r-1)+3},\quad
\lambda_{5,4r+3}^{4r+1}\lambda_{4p+1,4q+1}^{4r+1},\quad
\lambda_{5,4r+3}^{4r+3}\lambda_{4p+1,4q+1}^{4r+3}\\
&\lambda_{5,4r+3}^{4r+5}\lambda_{4p+1,4q+1}^{4(r+1)+1},\quad
\lambda_{5,4r+3}^{4r+7}\lambda_{4p+1,4q+1}^{4(r+1)+3}
\end{aligned}\right\}
$$
$$
\sim\left\{\begin{aligned}
 & \lambda_{5,4p+1}^{4p-3}\lambda_{4r+3,4q+1}^{4(p-1)+1},\quad
\lambda_{5,4p+1}^{4p-1}\lambda_{4r+3,4q+1}^{4(p-1)+3},\quad
\lambda_{5,4p+1}^{4p+1}\lambda_{4r+3,4q+1}^{4p+1}\\
&\lambda_{5,4p+1}^{4p+3}\lambda_{4r+3,4q+1}^{4p+3},\quad
\lambda_{5,4p+1}^{4p+5}\lambda_{4r+3,4q+1}^{4(p+1)+1}
\end{aligned}
\right\}.
$$

Since we have Lemma (\ref{5,i,i}), Lemma (\ref{5,i,i+4}) and assumption against previous Lemma,
we already know $\lambda_{5,i}^j$ for all $i,j$.
So we can simplify the equation as
$$\left\{\lambda_{4p+1,4q+1}^{4(r-1)+3},\quad
\lambda_{4p+1,4q+1}^{4(r+1)+3}
\right\}
\sim 
\left\{\lambda_{4r+3,4q+1}^{4(p-1)+1},\quad
\lambda_{4r+3,4q+1}^{4(p+1)+1}\right\}.$$

% Put $(p,q,r)=(p',q',r')$, and 
By assumption we get $\lambda_{4(p+1)+1,4q+1}^{4r+3}\sim\lambda_{4p+1,4q+1}^{4(r+1)+3}$.
%\sim(5,4p+1,4r+3,4q+1)$,
So we have
$$\lambda_{4(p+1)+1,4q+1}^{4r+3}\sim\lambda_{4p+1,4q+1}^{4(r+1)+3}\sim\cdots\sim\lambda_{5,4q+1}^{4p+4r+3}\sim 0,$$
the last relation coming from the previous lemma that
%  $N_{i,j}^3=0$ if $|i-j|>2$ and 
$\lambda_{5,i+2}^i=0$
and from fusion rules that $N_{i,j}^5=0$ if $|i-j|>4$.
% we have known 
% $\lambda_{a,b}^3$ and 
% $\lambda_{a,b}^5$ for all $a,b$.

Put $r=r'$ and we get what we need.
\end{proof}

% % Then consider $(3,4r+1,4p+1,4q+1)\sim(3,4p+1,4r+1,4q+1)$.
% % %\begin{equation}{\label{111}}
% % $$
% % \lambda_{4p+1,4q+1}^{4(r-1)+3}+\lambda_{4p+1,4q+1}^{4r+3}
% %     \sim
% %     \lambda_{4(p-1)+3,4q+1}^{4r+1}+\lambda_{4p+3,4q+1}^{4r+1}.
% % $$
% % %\end{equation}
% % We have $\lambda_{4p+3,4q+1}^{4r+1}=0$ for every $p\leq p',q\leq q',r\leq r'$.
% % Note that we can swap $p$ and $q$ as we need.

% Then consider $(3,4r+1,4p+1,4q+1)$ and $(3,4p+1,4q+1,4r+1)$,
% $$\lambda_{4p+1,4q+1}^{4(r-1)+3}+\lambda_{4p+1,4q+1}^{4r+3}
% \sim
% \lambda_{4(p-1)+3,4q+1}^{4r+1}+\lambda_{4p+3,4q+1}^{4r+1}
% %\lambda_{4(p-1)+3,4q+1}^{4r+1}+\lambda_{4p+3,4q+1}^{4r+1}
% \sim
% \lambda_{4p+1,4(q-1)+3}^{4r+1}+\lambda_{4p+1,4q+3}^{4r+1}.$$
% We have $\lambda_{4p+3,4q+1}^{4r+1}=0$ for every $p\leq p',q\leq q',r\leq r'$.


% Now consider $(3,4r+1,4p+1,4q+3)$ and $(3,4p+1,4q+1,4r+3)$,
% $$
% \begin{cases}
%     \lambda_{4p+1,4q+3}^{4(r-1)+3}+\lambda_{4p+1,4q+3}^{4r+3}
% \sim
% \lambda_{4(p-1)+3,4q+3}^{4r+1}+\lambda_{4p+3,4q+3}^{4r+1}\\
% \lambda_{4(p-1)+3,4q+3}^{4r+3}+\lambda_{4p+3,4q+3}^{4r+3}
% \sim
% \lambda_{4p+1,4(q-1)+3}^{4r+3}+\lambda_{4p+1,4q+3}^{4r+3}.
% \end{cases}
% %\lambda_{4(p-1)+3,4q+1}^{4r+1}+\lambda_{4p+3,4q+1}^{4r+1}
% $$
% %We have $\lambda_{4p+3,4q+1}^{4r+1}=0$ for every $p\leq p',q\leq q',r\leq r'$.

% % $(3,4p+1,4q+1,4r+3)$:
% % \lambda_{4(p-1)+3,4q+1}^{4r+3}+\lambda_{4p+3,4q+1}^{4r+3}
% % \sim
% % \lambda_{4p+1,4(q-1)+3}^{4r+3}+\lambda_{4p+1,4q+3}^{4r+3}.

% Consider $(3,4p+3,4q+1,4r+3)$ and $(3,4r+3,4p+3,4q+1)$,
% we have
% $$
% \begin{cases}
%     \lambda_{4p+1,4q+1}^{4r+3}+\lambda_{3,4p+3}^{4p+3}\lambda_{4p+3,4q+1}^{4r+3}+
%     \lambda_{4(p+1)+1,4q+1}^{4r+3}
%     \sim\lambda_{4p+3,4(q-1)+3}^{4r+3}+\lambda_{4p+3,4q+3}^{4r+3},\\
%     \lambda_{4p+3,4q+1}^{4r+1}+\lambda_{3,4r+3}^{4r+3}\lambda_{4p+3,4q+1}^{4r+3}+
%     \lambda_{4p+3,4q+1}^{4(r+1)+1}
%     \sim
%     \lambda_{4(p-1)+3,4q+1}^{4r+3}+\lambda_{4p+3,4q+1}^{4r+3}.
% \end{cases}$$

% If $\lambda_{3,3}^3= 0$,
% put $r\leq r'-1$ and we get
% $\lambda_{4p+3,4q+3}^{4r+3}=0$ and $\lambda_{4(p+1)+1,4q+1}^{4r+3}=0$ 
% for every $p\leq p'$, $q\leq q'$, and $r\leq r'-1$.


% % Consider $(5,4r'+3,4p'+1,4q'+3)$ and we get
% % $$\lambda_{4p'+1,4q'+3}^{4(r'-1)+3}+\lambda_{4p'+1,4q'+3}^{4(r'+1)+3}\sim
% % \lambda_{4(p'-1)+1,4q'+3}^{4r'+3}+\lambda_{4(p'+1)+1,4q'+3}^{4r'+3}.$$

% Consider $(3,4r+3,4p+1,4q+3)\sim(3,4p+1,4r+3,4q+3)$
% and we have
% \begin{equation}{\label{313}}
%     \lambda_{4p+1,4q+3}^{4r+1}+\lambda_{3,4r+3}^{4r+3}\lambda_{4p+1,4q+3}^{4r+3}+
%     \lambda_{4p+1,4q+3}^{4(r+1)+1}
%     \sim
%     \lambda_{4(p-1)+3,4q+3}^{4r+3}+\lambda_{4p+3,4q+3}^{4r+3}.
% \end{equation}
% If $\lambda_{3,3}^3= 0$,
% put $r\leq r'-1$ and we get
% $\lambda_{4p+3,4q+3}^{4r+3}=0$ for every $p\leq p',q\leq q',r\leq r'-1$.

% (-----------------------------)

% Now consider $(3,4r+3,4p+1,4q+1)\sim(3,4p+1,4r+3,4q+1)$.
% \begin{equation}{\label{311}}
%     \lambda_{4p+1,4q+1}^{4r+1}+\lambda_{3,4r+3}^{4r+3}\lambda_{4p+1,4q+1}^{4r+3}+
%     \lambda_{4p+1,4q+1}^{4(r+1)+1}
%     \sim
%     \lambda_{4(p-1)+3,4q+1}^{4r+3}+\lambda_{4p+3,4q+1}^{4r+3}.
% \end{equation}

%Now assume $\lambda_{4p+1,4q+1}^{4r+3}=0$ for $r\leq t$.
%When $r=t+1$,

% Assume $\lambda_{x,y}^{4r+1}=0$ if $|x-y|=4k+2$ for all $r\leq t$,
% then the right hand side of (\ref{111}) is $0$, thus $\lambda_{4p+1,4q+1}^{4r+3}=0$.

% Assume $\lambda_{x,y}^{4r-1}=0$ if $x\equiv y\equiv1 (\mod 4)$,
% then the right hand side of (\ref{111}) is $0$, thus $\lambda_{x,y}^{4r+3}=0$.
% \end{proof}

\begin{lemma}
$\lambda_{i,i}^3\neq 0$ for any odd $i$ compatible with fusion rules.
%\leq 2K+1$.
\end{lemma}
\begin{proof}
We already have $\lambda_{3,3}^3\sim\lambda_{5,5}^3\sim 1$.
Assume $\lambda_{k,k}^3\sim 1$ for $k\leq t$, where $t\geq 5$, but $\lambda_{t+2,t+2}^3=0$.

%Consider $(3,3,i,i+2)$ we get
Recall the first equation in (\ref{(3,3,i,i+2)}).
%    $$1+\lambda_{i,i+2}^5\sim
%    \lambda_{i,i}^3+\lambda_{i+2,i+2}^3.$$
Put $i=t$, and we get $\lambda_{t,t+2}^5=0$.
Put $i=t+2$, and we get $\lambda_{t+2,t+4}^5=0$.


Consider $(3,5,i,i)$, $\forall i\geq 3$,
%$$\lambda_{i,i}^3+\lambda_{i,i}^5+\lambda_{i,i}^7
%\sim\lambda_{5,i}^{i-2}+\lambda_{3,i}^i\lambda_{5,i}^i+\lambda_{5,i}^{i+2}.$$
$$\left\{1,\quad
\lambda_{i,i}^7\right\}
\sim
\left\{\lambda_{5,i}^{i-2},\quad
\lambda_{5,i}^{i+2}\right\}.
$$   
% we have
% $$\lambda_{i+2,i+2}^3+\lambda_{i+2,i+2}^5+\lambda_{3,5}^7\lambda_{i+2,i+2}^7
% \sim\lambda_{5,i+2}^i+
% \lambda_{3,i+2}^{i+2}\lambda_{5,i+2}^{i+2}+
% \lambda_{5,i+2}^{i+4}.$$


%Since $\lambda_{5,5}^3\sim 1$, 
Put $i=t+2$ and we have
$\{1,\lambda_{t+2,t+2}^7\}\sim\{0,0\}.$
We get the contradiction by Theorem (\ref{samerank}).
\end{proof}

\begin{theorem}
If $m$ is even, then $\lambda_{i,j}^k\neq 0$ for every $N_{i,j}^k =1$.
\end{theorem}
\begin{proof}
The previous lemma asserts Lemma (\ref{recursion}). 
We can get what we want by recursion.
\end{proof}

\section{The uniqueness of $(\mathcal{U},Y)$}
In this section, we fix basis $\{e_{i,j}\}$ and 
$Y(\cdot,z)$ a VOA structure on $\mathcal{U}$ such that $\lambda_{i,j}^k=1$
for all $i,j,k$ compatible with fusion rules.
Assume there is another VOA structure $\overline{Y}(\cdot,z)$ on $\mathcal{U}$,
and we will prove that the square of $\lambda_{b,c}^a$ for $\overline{Y}(\cdot,z)$
equals $1$ whenever $N_{b,c}^a=1$
if braiding matrix satisfies some non-vanishing conditions.
These conditions can be checked case by case for any fixed $\mathcal{U}_k$,
and we get $(\mathcal{U},\overline{Y}(\cdot,z))\cong (\mathcal{U},Y(\cdot,z))$,
thus proving the uniqueness of $\mathcal{U}$.

\begin{lemma}
$\lambda_{1,k}^k=\lambda_{k,1}^k=\lambda_{k,k}^1=1$ for all $k$.
\end{lemma}
\begin{proof}    
Since we assume $\mathcal{U}_{n-1}$ has a unique VOA structure,
we have $Y(u,z)=\overline{Y}(u,z)$ for all $u\in U^1$, that is $\lambda_{1,k}^k=1$.    

For any $u^{1}\in U^{1}$, $u^{k}\in U^{k}$, skew symmetry of $Y\left(\cdot,z\right)$
and $\overline{Y}\left(\cdot,z\right)$ ( \cite{FHL} ) imply

$$
\overline{Y}(u^{k},z)u^{1}=e^{zL\left(-1\right)}\overline{Y}\left(u^{1},-z\right)u^{k}=e^{zL\left(-1\right)}Y(u^{1},-z)u^{k}=Y\left(u^{k},z\right)u^{1}=\mathcal{I}_{k,1}^{k}\left(u^{k},z\right)u^{1}.
$$
In the mean time, $\overline{Y}\left(u^{k},z\right)u^{1}=\lambda_{k,1}^{k}\cdot\mathcal{I}_{k,1}^{k}(u^{k},z)u^{1}$.
Thus we get $\lambda_{k,1}^{k}=1$.

Note that by Remark \ref{self-dual}, $\mathcal{U}$ has a unique
invariant bilinear form $\left\langle \cdot,\cdot\right\rangle $
with $\left\langle 1,1\right\rangle =1$. For $u^{1}\in U^{1}$ and
$u^{k},v^{k}\in U^{k}$, we have
$$
\left\langle Y\left(u^{k},z)v^{k}\right),u^{1}\right\rangle =\left\langle v^{k},Y\left(e^{zL\left(-1\right)}\left(-z^{-2}\right)^{L\left(0\right)}u^{k},z^{-1}\right)u^{1}\right\rangle .
$$
That is,
$$
\left\langle \mathcal{I}_{k,k}^{1}\left(u^{k},z\right)v^{k},u^{1}\right\rangle 
=\left\langle v^{k},\mathcal{I}_{k,1}^{k}\left(e^{zL\left(-1\right)}\left(-z^{-2}\right)^{L\left(0\right)}u^{k},z^{-1}\right)u^{1}\right\rangle .
$$
The invariant bilinear form on $\left(\mathcal{U},\overline{Y}\right)$
gives
$$
\left\langle \lambda_{k,k}^{1}\cdot\mathcal{I}_{k,k}^{1}\left(u^{k},z\right)v^{k},u^{1}\right\rangle 
=\left\langle v^{k},\lambda_{k,1}^{k}\cdot\mathcal{I}_{k,1}^{k}\left(e^{zL\left(-1\right)}\left(-z^{-2}\right)^{L\left(0\right)}u^{k},z^{-1}\right)u^{1}\right\rangle .
$$
Since $\lambda_{k,1}^k=1$, we get $\lambda_{k,k}^{1}=1$.
\end{proof}

\begin{theorem}
$\lambda_{i,j}^k=\lambda_{j,i}^k=\lambda_{j,k}^i$ for all $i,j,k$ compatible with fusion rules.
\end{theorem}
\begin{proof}
This comes from the proof of Theorem (\ref{alternation}) and previous lemma.
\end{proof}

\begin{lemma}\label{partunitary}
$B^{\operatorname{T}}\tilde{B}=I$.
\end{lemma}
\begin{proof}
Recall (\ref{braid matrix and lambda}).
This is the direct consequence from the $Y(\cdot,z)$ we chose.
\end{proof}

% \begin{lemma}
%     $\left(\lambda_{a,b}^{\beta_{i}}\right)^2=\left(\lambda_{a,b}^{\mu_i}\right)^2$ for all $i=1,\ldots,N$.
% \end{lemma}
% \begin{proof}
%     Consider the case $a=d$ and $b=c$, 
%     and we get
%     $$B\operatorname{diag}\left(\left(\lambda_{b,\beta_1}^a\right)^2,\ldots,\left(\lambda_{b,\beta_N}^a\right)^2\right)
% \tilde{B}^{\operatorname{T}}=\operatorname{diag}\left(\left(\lambda_{b,\mu_1}^a\right)^2,\ldots,\left(\lambda_{b,\mu_N}^a\right)^2\right).$$

%     By previous lemma,
%     $\left(\lambda_{a,b}^{\beta_{i}}\right)^2=\left(\lambda_{a,b}^{\mu_i}\right)^2$ for all $i=1,\ldots,N$.
% \end{proof}


Fusion rules show that the longest direct summand appears at the $(2K+1)$st position when $m=4K+1$ and at the $(2K+3)$rd position when $m=4K+3$.
When $m=4K$, the longest direct summand appears at the $(2K+1)$st position. 
When $m=4K+2$, the longest direct summand appear at both the $(2K+1)$st and $(2K+3)$rd position,
and in this paper we use the $(2K+1)$st one.
We denote the module with the longest fusion rule by $U^t$.

% \begin{proposition}
%     $B_{U^t,U^1}\neq 0.$
% \end{proposition}

\begin{lemma}
$\left(\lambda_{t,t}^a\right)^2=1$ if $\left(\tilde{B}_{a,t}^{a,t}\right)_{U^1,U^t}\neq 0$.
\end{lemma}
\begin{proof}
Consider $(a,b,t,t)$ in (\ref{braid matrix and lambda}), and we have
$$
\operatorname{diag}(\lambda_{a,b}^{\beta_1}\lambda_{t,t}^{\beta_1},\ldots,\lambda_{a,b}^{\beta_N}\lambda_{t,t}^{\beta_N})\tilde{B}
=\tilde{B}\operatorname{diag}(\lambda_{a,t}^{\mu_1}\lambda_{b,t}^{\mu_1},\ldots,\lambda_{a,t}^{\mu_N}\lambda_{b,t}^{\mu_N}).$$
Compare the $(i,j)$-element in the matrix of each side of the euqation,
and we have 
\begin{equation}\label{abtt}
\lambda_{a,b}^{\beta_i}\lambda_{t,t}^{\beta_i}\tilde{B}_{i,j}=\tilde{B}_{i,j}\lambda_{a,t}^{\mu_j}\lambda_{b,t}^{\mu_j}
\end{equation}
for all $i$ and $j$.
Choose $(i,j)$ such that $\beta_i=U^1$ and $\mu_j=U^t$.
It is possible by computing the fusion rules.
% noting that when $m=4k+2$, it will be two possible $t$ and we can choose either we want.
Since $\lambda_{t,t}^1=1$ and $b=a$, 
we get $\tilde{B}_{i,j}=\tilde{B}_{i,j}\left(\lambda_{a,t}^{\mu_i}\right)^2$.
So we have $\left(\lambda_{t,t}^a\right)^2=1$ if 
$\left(\tilde{B}_{a,t}^{a,t}\right)_{U^1,U^t}\neq 0$. 
\end{proof}

\begin{theorem}\label{squarelambda}
For all $i,j,k$ compatible with fusion rules, 
$\left(\lambda_{p,q}^r\right)^2=1$ if 
$\left(\tilde{B}_{p,t}^{q,t}\right)_{U^r,U^t}\neq 0$.
\end{theorem}
\begin{proof}
Choose $j$ such that $\mu_j=U^t$ in (\ref{abtt}),
and go through all $\beta_i$ such that $N_{a,b}^{\beta_i}=1$.
We get what we want by multiplying $\lambda_{t,t}^{\beta_i}$ on both sides of the equation
and setting $(p,q,r)=(a,b,\beta_i)$.
\end{proof}

\begin{remark}\label{rmkB}
% Since $B=\left(\tilde{B}^{\operatorname{T}}\right)^{-1}$ by Lemma (\ref{partunitary}),
% We can always use elements in $\tilde{B}$ instead of $B$.
From now on if we use 
%elements in $B$, or equaly by (\ref{braid matrix and lambda}), 
elements in $\tilde{B}$, we will assume they are non-zero.
For any fixed $\mathcal{U}_k$, it can be checked case by case. 
Examples of computing elements of braiding matrix are given at the end of this section.

% Samely, we can also just check elements in $\tilde{B}^{\operatorname{T}}$.
\end{remark}

%We need to show it is well defined.
%If $N_{a,b}^k=1$, and $N_{a',b'}^k=1$, we need to show
%$\lambda_{}$

%If $m\neq 4k+2$

\begin{lemma}
Define a map
$\sigma:(\mathcal{U},Y)\to(\mathcal{U},\overline{Y})$ 
by
\begin{itemize}
\item If $N_{t,t}^k=1$, then $\sigma(u^k):=\lambda_{t,t}^k u^k$;
\item If $N_{t,t}^k=0$, then $\sigma(u^k):=\lambda_{a,b}^k \sigma(u^a)_{(n)}\sigma(u^b)\restriction U^k$ for any $N_{a,b}^k=1$ and $u^a_{(n)}u^b\restriction U^k=u^k$ (Here $\restriction$ is the restriction),
\end{itemize} 
and linear expansion.
Then it is a well-defined homomorphism.
\end{lemma}
\begin{proof}
We write $u^p_{(n)}u^q=\sum_r w^r$,
and we write the $n$th product of $\overline{Y}$ as $(\overline{n})$.
By the definition of $\lambda_{p,q}^r$ we have $u^p_{(\overline{n})}u^q=\sum_r \lambda_{p,q}^r w^r$.
Note that from (\ref{abtt}) and by putting $\mu_j=t$,
we always have 
\begin{equation}\label{lambdat}
\lambda_{a,b}^{\beta_i}\lambda_{t,t}^{\beta_i}=\lambda_{a,t}^{t}\lambda_{b,t}^{t}.
\end{equation}

% For all $u\in U^a,v\in U^b$ and $w=\sigma(u_{(n)}v)\restriction U^k$ such that $N_{a,b}^k=1$, we have
% $\sigma(u_{(n)}v)\restriction U^k=\lambda_{t,t}^k w$
% and $\sigma(u)_{(n)}\sigma(v)\restriction U^k
% =\lambda_{t,t}^a\lambda_{t,t}^b u_{(n)}v\restriction U^k
% =\lambda_{t,t}^a\lambda_{t,t}^b w$,
% thus $\lambda_{t,t}^k=\lambda_{t,t}^a\lambda_{t,t}^b$.

If $m\neq 4K+2$, then it is well-defined since the fusion rule that $N_{t,t}^k=1$ for all odd $k$.
If $m=4K+2$, then $t=2K+1$ and $k=m+1$ is the only case when $N_{t,t}^k=0$.
Assume there are $a',b'<k$ such that $N_{a',b'}^k=1$ and $u^{a'}_{(n)}u^{b'}\restriction U^k=u^k$.
We need to prove $\lambda_{a,b}^k\lambda_{t,t}^a\lambda_{t,t}^b =\lambda_{a',b'}^k\lambda_{t,t}^{a'}\lambda_{t,t}^{b'},$
and it is obvious by (\ref{lambdat}). So we get the well-defined map.

Now we check $\sigma(u^p_{(n)}u^q)=\sigma(u^p)_{(\overline{n})}\sigma(u^q)$.
We have $\sigma(u^p_{(n)}u^q)=\sum_r \lambda_{t,t}^r w^r$, 
and $\sigma(u^p)_{(\overline{n})}\sigma(u^q)=\lambda_{t,t}^p\lambda_{t,t}^q\sum_r \lambda_{p,q}^r w^r$.
We get the equation by (\ref{lambdat}) and Theorem (\ref{squarelambda}).

% Consider $(a',b',a,b)\sim(a',a,b',b)$ and Theorem (\ref{samerank}), 
% there exists $\beta$ such that 
% $\lambda_{a',b'}^{k}\lambda_{a,b}^{k}=\lambda_{a,a'}^\beta\lambda_{b,b'}^\beta$.
%Thus $\lambda_{t,t}^\beta=\lambda_{t,t}^{\mu}\lambda_{t,t}^{\mu'}$ for $\mu=a,b$.
%Since $\lambda_{p,q}^r\neq 0$ if $N_{p,q}^r\neq 0$, 
%we get what we want 
%The proof is complete by multiplying $\lambda_{a,b}^k\lambda_{t,t}^{a'}\lambda_{t,t}^{b'}$
%on each side of the equation above.
%Square both sides of the equation, 
%and we have $\left(\lambda_{a',b'}^{k}\right)^2=\left(\lambda_{a,b}^{m+1}\right)^2$
%by previous lemma.
\end{proof}

\begin{theorem}\label{unique}
The VOA structure of $\mathcal{U}$ is unique 
if corresponding elements in braiding matrices are non-zero.
\end{theorem}
\begin{proof}
%From (\ref{squarelambda}) and previous lemma, we found $\sigma$ such that $\sigma^2=\operatorname{id}$.
The $\sigma$ from the previous lemma is the isomorphism we want.
It is the direct consequence of previous lemma and Theorem (\ref{squarelambda}).
% Since $\sigma$ is a linear map,
% we only need to check the case when $u^a_{(n)}u^b=u^k$.
% We have
% $$\sigma(u^a)_{(n)}\sigma(v^b)=\lambda_{t,t}^a\lambda_{t,t}^b u^a_{(n)}v^b=\lambda_{t,t}^a\lambda_{t,t}^b\lambda_{a,b}^k u^k$$
% and
% $$\sigma(u^a_{(n)}v^b)=\lambda_{a,b}^{k}\sigma(u^k)=\lambda_{a,b}^{k}\lambda_{t,t}^{k}u^k.$$
% Since $\lambda_{t,t}^k=\lambda_{t,t}^a\lambda_{t,t}^b$, the proof is complete.
% This makes $\overline{Y}=Y^{\sigma}$.
% By (\ref{new structure}) we get $(\mathcal{U},Y)\cong(\mathcal{U},\overline{Y})$.
\end{proof}

Note that we have (\ref{abtt}) and Remark (\ref{rmkB}), 
so it is enough to check elements of $\tilde{B}$ are non-zero.
We state examples of computing $\tilde{B}$, 
especially when $\mathcal{U}=\mathcal{U}_1=\mathcal{U}_{\operatorname{6A}}$,
and it can be found in \cite{DJY}. And we can get the uniqueness of $\mathcal{U}_2$ as an example
by direct computation following Proposition (\ref{fusion rules of virasoro modules}) and (\ref{FFK 2.19}).

\begin{example}
Consider braiding matrix for $L\left(25/28,0\right)$-modules.
Denote irreducible $L\left(25/28,0\right)$-modules $L\left(25/28,34/7\right)$
and $L\left(25/28,9/7\right)$ by $Q_{2}$ and $Q_{3}$
respectively. For convenience, we will denote $\left(\tilde{B}_{Q_{a},Q_{b}}^{Q_{c},Q_{d}}\right)_{Q_{e},Q_{f}}$
by $\left(\tilde{B}_{a,b}^{c,d}\right)_{e,f}$, $a,b,c$, $d,$ $e$,
$f$ $\in\left\{ 2,3\right\} $. 
Then we have
\label{(2332)_3,2 nonzero} \label{(3322)_3,3 nonzero}
\label{(2333)_3,2 nonzero} \label{(3,2,3,3)_2,3 nonzero}$\left(\tilde{B}_{2,2}^{3,3}\right)_{3,2}\not=0$,
$\left(\tilde{B}_{3,2}^{3,2}\right)_{3,3}\not=0$, $\left(\tilde{B}_{2,3}^{3,3}\right)_{3,2}\not=0$,
and $\left(\tilde{B}_{3,3}^{2,3}\right)_{2,3}\not=0$.
\end{example}

\begin{example}
We use notations in Remark (\ref{rmk}), and consider $L(11/12, 0)$-modules.
In this case $t=3$,
and we have
$\left(\tilde{B}_{a,3}^{b,3}\right)_{3,c}\neq 0$
for all $a,b,c\in \{3,5,7\}$ such that those $\tilde{B}_{i,j}$ 
appear in Theorem (\ref{squarelambda}) and compatible with fusion rules.


We write the braiding matrix as $B_{(a,b,c)}$.
We need to check all the 
$$(a,b,c)\in\left\{\begin{aligned}
&(3,3,3),(3,3,5),(3,5,3),(3,5,5),\\
&(5,3,3),(5,3,5),(5,5,3),(5,5,5),\\
&(5,5,7),(5,7,5),(7,5,5)
\end{aligned}\right\}.$$

To prove $B_{(3,3,3)}=\left(\tilde{B}_{3,3}^{3,3}\right)_{3,3}\neq 0$, 
we need to check 
$$
r(3,3,3,3)_{3,3}=r(3,2,3,2)_{2,3}r(2,2,3,3)_{3,2}+r(3,2,3,4)_{2,3}r(2,2,3,3)_{3,4}\neq 0.
$$
Since 
$$\begin{aligned}
r(3,2,3,2)_{2,3}&=r(3,2,2,1)_{2,2}r(2,2,2,2)_{1,3}=x^{1/4}x^{-1/4}\frac{[1]}{[2]},\\
r(2,2,3,3)_{3,2}&=r(2,2,2,2)_{3,1}r(1,2,2,3)_{2,2}+r(2,2,2,2)_{3,3}r(3,2,2,3)_{2,2},\\
&=x^{-1/4}\frac{[3]}{[2]}x^{1/4}+\left(-x^{-1/4-2/2}\frac{[1]}{[2]}x^{-1/4+3/2}\frac{[1]}{[3]}\right),\\
r(3,2,3,4)_{2,3}&=r(3,2,2,3)_{2,2}r(2,2,2,4)_{3,3}+r(3,2,2,3)_{2,4}r(4,2,2,4)_{3,3}\\
&=x^{-1/4+3/2}\frac{[1]}{[3]}x^{1/4}+x^{-1/4}\frac{[2]}{[3]}x^{-1/4+4/2}\frac{[1]}{[4]},\\
r(2,2,3,3)_{3,4}&=r(2,2,2,2)_{3,3}r(3,2,2,3)_{2,4}\\
&=-x^{-1/4-2/2}\frac{[1]}{[2]}x^{-1/4}\frac{[2]}{[3]},
\end{aligned}$$
let $r=x^{1/2}+x^{-1/2}$, and we get 
$$
\begin{aligned}
%\left(\tilde{B}_{3,3}^{3,3}\right)_{3,3}
B_{(3,3,3)}
&=\frac{[1]}{[2]}\left(\frac{[3]}{[2]}-\frac{[1][1]}{[2][3]}\right)-\left(\frac{[1]}{[3]}+\frac{[1][2]}{[3][4]}\right)\frac{[1]}{[3]}\\
&=1-\frac{1}{r^2-1}-\frac{1}{(r^2-1)^2}-\frac{1}{(r^2-1)^2(r^2-2)}\\
&=1-\frac{1}{r^2-2}.
\end{aligned}
$$
That is, $\operatorname{Re}B_{(3,3,3)}> 0$, so $B_{(3,3,3)}\neq 0$

Following the same approach above, we have
$$\begin{aligned}
\operatorname{Im}B_{(3,3,5)}>0,&
\operatorname{Re}B_{(3,5,3)}<0,&
\operatorname{Im}B_{(3,5,5)}>0,&
\operatorname{Im}B_{(5,3,3)}>0,&
\operatorname{Re}B_{(5,3,5)}>0,\\
\operatorname{Im}B_{(5,5,3)}>0,&
\operatorname{Re}B_{(5,5,5)}<0,&
\operatorname{Re}B_{(5,5,7)}>0,&
\operatorname{Re}B_{(5,7,5)}<0,&
\operatorname{Re}B_{(7,5,5)}<0.
\end{aligned}$$

% The case 
% $(3,3,3),(3,5,3),(5,3,5),(5,5,5),(5,5,7),(5,7,5)$ 
% can be checked since their real parts are non-zero, 
% and other cases have non-zero imaginary parts.

So all the elements in braiding matrix we need are non-zero. 
Thus $\mathcal{U}_2=U_1^1\oplus U_1^3\oplus U_1^5\oplus U_1^7=
\left[0,0\right]_1\oplus\left[2,\frac{5}{4}\right]_1
\oplus\left[4,\frac{19}{4}\right]_1\oplus\left[6,\frac{21}{2}\right]_1$
has a unique VOA structure
by Theorem (\ref{squarelambda}) and the previous theorem.
\end{example}

\section{Generated by Griess algebra}
It is proved in \cite{JZ} that $\mathcal{U}_1$ is generated by Griess algebra.
Now we prove that if $\mathcal{U}_{k-1}$ is generated by Griess algebra, so is $\mathcal{U}=\mathcal{U}_k$.
\begin{lemma}
$U^3_{k-1}$ is generated by Griess algebra.
\end{lemma}
\begin{proof}
%    Recall notations in (\ref{rmk}). 
% Since $U^3=P^3\otimes Q^3$ and $Q^3$ is generated by the Virasoro vector,
% it's suffice to show $P^3$ is generated by Griess algebra.
Since $U^3_{k-1}$ is an irreducible $\mathcal{U}_{k-1}$-module,
we only need to find one vector with weight $2$ in it.

For $\mathcal{U}_1$, 
we have $U^{3}_0=\mathcal{V}(\frac{5}{7})\otimes L(\frac{25}{28},\frac{9}{7})$,
and $\frac{5}{7}+\frac{9}{7}=2.$

For $\mathcal{U}_k$ with $k\geq 2$, we have
$U^3_{k-1}=\mathcal{U}_{k-1,2}\otimes L\left(c_{k+6},h_{(1,3)}^{k+6}\right)$
and 
$$\mathcal{U}_{k-2,0}
\otimes L\left(c_{k+5},h_{(3,1)}^{k+5}\right)
=\left[0,h_{(3,1)}^{k+5}\right]_{k-2}\subseteq \mathcal{U}_{k-1,2}.$$
%     From GKO construction we know that 
% % $$\mathcal{U}_{1,2}\otimes\left(\bigotimes_{i=3}^{k} L\left(k+7,h^{(k+7)}_{(i,1)}\right)\right)\subseteq U_k^3.$$
% $$\mathcal{V}\otimes L\left(\frac{25}{28},\frac{3}{4}\right)\otimes\left(\bigotimes_{i=3}^{k} L\left(k+7,h^{(k+7)}_{(i,1)}\right)\right)\subseteq U_k^3.$$
Note that
$$
h^{(k+6)}_{(1,3)}+ h^{(k+5)}_{(3,1)}=\frac{k+8}{k+6}+\frac{k+4}{k+6}=2,
$$
and the proof is complete.
\end{proof}

\begin{theorem}
$\mathcal{U}$ is generated by Griess algebra.
\end{theorem}
\begin{proof}
We already know $U^1$ and $U^3$ are generated by Griess algebra.
Assume $U^{2k+1}$ are generated by Griess algebra for $0\leq k\leq n$.
By fusion rules of $U^3$ we have $U^{2k+3}\subseteq U^3\boxtimes U^{2k+1}$,
so $U^{2n+3}$ is generated by Griess algebra.
Thus $\mathcal{U}$ is generated by Griess algebra.
\end{proof}

\section{Open problems}
Here are some open problems based on 
this paper and its citations.

\begin{problem}
	Quantum dimension theory \cite{DJX} is used
	in \cite{DJY} to determine irreducible modules and fusion rules of $\mathcal{U}_{6A}$. 
	For this paper, is there a good way to determine irreducible modules and fusion rules of $\mathcal{U}_k$ for all $k$ simultaneously?
\end{problem}

\begin{problem}
	So far, all elements of braiding matrices we computed are nonzero, 
	is this a general result?
\end{problem}

\begin{problem}
	It is proved in \cite{LYY, JZ} that $\mathcal{U}_1$ is generated by Ising vectors while
	we only proved that $\mathcal{U}_k$ are generated by Griess algebra. 
	Are all $\mathcal{U}_k$ generated by Ising vectors?
	The answer for this problem can help proving Conjecture A.14 in \cite{LY}.
\end{problem}

\section*{Acknowledgments}
W. Zheng is supported by the NSFC No.~12201334 and the National Science Foundation of Shandong Province No.~ZR2022QA023.
% Note that $U^{m+1}$ will be 
% simple current modules if $m$ is even,
% and this make the proof miserable in the last part of Section 4.
% It might exist a simplified proof of it since its quantum dimension is equal to $1$.

% Most steps in this paper
% does not involve detailed structures of 
% vertex operator algebras or their modules 
% as long as it has a unique VOA structure 
% and can provide enough modules to pair with
% part of unitary series.
% And \cite{Lin} gives the powerful theorem (\ref{c1}),
% with which we only need subsititutions to prove (\ref{same fusion rule}).

% It makes sense if such uniqueness theorems
% could be extended to other VOA series 
% other than what we have constructed. 
% Furthermore, it is predicted that 
% there will be a categorical version of prove
% since the braiding matrices form a ribbon structure
% on subcategory of certain VOA modules.
% And these results have potential for classification of
% modular tensor categories \cite{R}.

\begin{thebibliography}{99}
\bibitem[ABD]{ABD} T. Abe, G. Buhl, C. Dong, Rationality, regularity
and $C_{2}$-cofiniteness. \emph{Trans. AMS.} \textbf{356} (2004),
3391\textendash 3402.

\bibitem[ADL]{ADL} T. Abe, C. Dong, H. Li, Fusion rules for the vertex
operator algebras $M(1)^{+}$ and $V_{L}^{+}$ . \emph{Comm. Math.
	Phys.} \textbf{253} (2005), 171\textendash 219.

\bibitem[ADJR]{ADJR} C. Ai, C. Dong, X. Jiao, L. Ren, The irreducible
modules and fusion rules for the parafermion vertex operator algebras.
\emph{Trans. Amer. Math. Soc. }\textbf{370} (2018), no. 8, 5963\textendash 5981.

%\bibitem[BK]{BK} B. Bakalov and A. Kirillov, Lectures on tensor categories and modular functors. {\em University Lecture Series} Vol 21, 2000.

\bibitem[B]{B} R.E. Borcherds,
Vertex algebras, Kac-Moody algebras and the Monster, 
Proc.Nat.Acad.Sc. \textbf{83} (1986) 306\textendash 3071.

\bibitem[C]{C} J. H. Conway, A simple construction for the Fisher-Griess
Monster group. \emph{Invent. Math. }\textbf{79 }(1985), 513\textendash 540.

%\bibitem[CKM]{CKM} T. Creutzig, S. Kanade and R. McRae, 
%Gluing vertex algebras,  \emph{Adv. Math.} \textbf{396} (2022), 108174, 72 pp.

\bibitem[DJX1]{DJX} C. Dong, X. Jiao, F. Xu, Quantum dimensions and
quantum Galois theory. \emph{Trans. Amer. Math. Soc.}\textbf{ 365}
(2013), 6441\textendash 6469.

\bibitem[DJX2]{DJX2} C. Dong, X. Jiao and F. Xu, Mirror extensions of vertex operator algebras, {\em Comm. Math. Phys. } {\bf329} (2014), 263-294.

\bibitem[DJY]{DJY} C. Dong, X. Jiao, N. Yu, 6A-algebra and its representations.
\emph{J. Algebra.}\textbf{533}(2019), 174\textendash 210.

\bibitem[DL]{DL} C. Dong, J. Lepowsky, Generalized vertex operator
algberas and relative vertex operators. \emph{Progress in Mathematics,
}vol. 112. Birkhauser Boston. Massachusetts, 1993.

\bibitem[DLM1]{DLM1} C. Dong, H. Li, G. Mason, Regularity of rational
vertex operator algebras. \emph{Adv. Math.} \textbf{132} (1997), 148\textendash 166.

\bibitem[DLM2]{DLM2} C. Dong, H. Li, G. Mason, Twisted representations
of vertex operator algebras. \emph{Math. Ann. }\textbf{310 }(1998),\textbf{
}571\textendash 600.

\bibitem[DLM3]{DLM3}C. Dong, H. Li, G. Mason, Simple currents and
extensions of vertex operator algebras. \emph{Comm. Math. Phys.} \textbf{180
}(1996), 671-707.

\bibitem[DLM4]{DLM4}C. Dong, H. Li, G. Mason, Compact automorphism groups of vertex operator algebras.
\emph{Internat. Math. Res. Notices.} \textbf{18
}(1996), 913-921.

\bibitem[DM]{DM} C. Dong, G. Mason, On quantum Galois theory.\emph{
Duke Math. J. }\textbf{86} (1997), 305\textendash 321.

\bibitem[DMZ]{DMZ} C. Dong, G. Mason, Y. Zhu, Discrete series of
the Virasoro algebra and the moonshine module. Proc. \emph{Symp. Pure.
Math., American Math. Soc.} \textbf{56} II (1994), 295\textendash 316.

%\bibitem[DNR]{DNR} C. Dong, S. Ng, L. Ren, Orbifolds and minimal modular extensions.
%\emph{Adv. Math.}\textbf{462} (2025), 110103.

\bibitem[DZ]{DZ} C. Dong, W. Zheng, Uniqueness of VOA structure of $3C$-algebra and
$5A$-algebra. \emph{J. Algebra
}\textbf{572 }(2021), 76\textendash 110.

\bibitem[FFK]{FFK} G. Felder, J. Frohlich, G. Keller, Braid matrices
and structure constants for minimal conformal models. \emph{Comm.
Math. Phys.} \textbf{124} (1989), no. 4, 647\textendash 664.

\bibitem[FHL]{FHL} I. B. Frenkel, Y. Huang, J. Lepowsky, On axiomatic
approaches to vertex operator algebras and modules. \emph{Memoirs
American Math. Soc. }\textbf{104}, 1993.

\bibitem[FLM]{FLM} I. B. Frenkel, J. Lepowsky, A. Meurman, Vertex
operator algebras and the monster. \emph{Pure and Applied Math., }vol.
134, Academic Press, Massachusetts, 1988.

\bibitem[FZ]{FZ} I. B. Frenkel, Y. Zhu, Vertex operator algebras
associated to representations of affine and Virasoro algebras. \emph{Duke
Math. J. }\textbf{66 }(1992), no.1, 123\textendash 168.

\bibitem[GKO]{GKO} P. Goddard, A. Kent, D. Olive, Unitary representations
of the Virasoro and super Virasoro algebras. \emph{Comm. Math. Phys.
}\textbf{103 }(1986), 105\textendash 119.

%\bibitem[H1]{H1} Y. Huang, A theory of tensor products for module
%categories for a vertex operator algebra, IV. \emph{J. Pure Appl.
%Algebra} \textbf{100 }(1995), 173\textendash 216.

%\bibitem[H2]{H2} Y. Huang, Virasoro vertex operator algebras, the
%(nonmeromorphic) operator product expansion and the tensor product
%theory. \emph{J. Algebra} \textbf{182} (1996), 201\textendash 234.

\bibitem[H1]{H1} Y. Huang, Generalized rationality and a ``jacobi
identity\textquotedblright{} for intertwining operator algebras. \emph{Selecta
Math.}\textbf{ 6 }(2000), 225\textendash 267.

\bibitem[H2]{H2} Y. Huang, Differential equations and intertwining operators. \emph{Commun.
	Contemp. Math.}\textbf{ 7 }(2005), 375\textendash 400.

\bibitem[HKL]{HKL} Y. Huang, A. Kirillov, J. Lepowsky, Braided tensor
categories and extensions of vertex operator algebras\emph{. Comm.
Math. Phys}. \textbf{337} (2015), no. 3, 1143\textendash 1159.

\bibitem[JZ]{JZ} X. Jiao, W. Zheng, Vertex operator algebras generated by two Ising vectors,
\emph{J. Algebra} \textbf{610} (2022), 546\textendash 570.

%\bibitem[KO]{KO} A. Kirillov and V. Ostrik, On a q-analog of Mckay correspondence and ADE classification of $\hat{sl}_{(2)}$ conformal field theories, {\em Adv. Math.} {\bf 171}(2002), 183-227.

\bibitem[KR]{KR} V.G. Kac, A.K. Raina, Bombay lectures on highest
weight representations of infinite dimensional Lie algebras. World
Scientific, Singapore, 1987.

\bibitem[KZ]{KZ} V.G. Knizhnik, A.B. Zamolodchikov, Current algebras
and wess-zumino model in two dimensions.\emph{ Nucl. Phys. B.} \textbf{247}
(1984), 83\textendash 103.

\bibitem[Li1]{L}  H. Li, Some finitenss properties of regular vertex operator algebras, {\em J. Algebra} {\bf 212} (1999), 495-514.

\bibitem[Li2]{Li}H. Li, Symmetric invariant bilinear forms on vertex
operator algebras. \emph{J. Pure Appl. Algebra} \textbf{96 }(1994),
279\textendash 297.

\bibitem[Lin]{Lin}X. Lin, Mirror extensions of rational vertex operator algebras.
\emph{Trans. Amer. Math. Soc.} \textbf{369} (2017), 3821\textendash 3840.

%\bibitem[LY1]{LY}  C.H. Lam and H. Yamada, Sigma involutions associated with parafermion vertex operator algebra
%$K(sl_2, k)$, \emph{Linear and Multilinear Algebra}, \textbf{70 (21)} (2022), 6780-6819.

\bibitem[LY]{LY} C.H. Lam and H. Yamauchi, 
The Conway-Miyamoto correspondences for the Fischer 3-transposition groups.
\emph{. Trans. Amer. Math. Soc. }\textbf{375} (2022), no. 3, 2025\textendash 2067.

\bibitem[LYY]{LYY} C. H. Lam, H. Yamada, H. Yamauchi, Vertex operator
algebras, extended $E_{8}$ diagram, and McKay's observation on the
monster simple group\emph{. Trans. Amer. Math. Soc. }\textbf{359}
(2007), no. 9, 4107\textendash 4123.

\bibitem[M]{M} M. Miyamoto, Vertex operator algebras generated by
two conformal vectors whose $\tau$-involutions generate $S_3$.
\emph{J. Algebra }\textbf{268 }(2003), 653\textendash 671.

%\bibitem[M]{M2} M. Muger, From subfactors to categories and topology. II. The quantum double of tensor categories and subfactors, {\em J. Pure Appl. Algebra} {\bf 180} (2003), 159-219.

%\bibitem[O]{O} V. Ostrik, Module categories, weak hopf algebras and modular invariants, {\em Transform. Groups} {\bf 8} (2003), 177-206.

%\bibitem[R]{R} B.C. Rayhaun; Bosonic rational conformal field theories in small genera, chiral fermionization, and symmetry/subalgebra duality. 
%\emph{J. Math. Phys.} \textbf{65} (2024), 052301.

\bibitem[SY]{SY} S. Sakuma, H. Yamauchi, Vertex operator algebra
with two Miyamoto involutions generating $S_{3}$. \emph{J. Algebra
}\textbf{267 }(2003), 272\textendash 297.

\bibitem[TK]{TK} A. Tsuchiya, Y. Kanie, Vertex operators in conformal
field theory on P1 and monodromy representations of braid group. \emph{Adv.
Stud. Pure Math.} \textbf{16} (1988) 297\textendash 372.

\bibitem[W]{W} W. Wang, Rationality of Virasoro vertex operator algebras\emph{.
Internat. Math. Res. Notices. }Yale University. Connecticut, 1990\emph{. }

\bibitem[X]{X1} F. Xu, Mirror extensions of local nets, {\em Commun. Math. Phys.} {\bf 270} (2007), 835-847.

\bibitem[Zheng]{Zheng} W. Zheng,
The VOAs generated by two Ising vectors e and f with $\langle e,f\rangle=\frac{1}{2^8}$,
$\tau_e\tau_f=3$ or $\langle e,f\rangle=\frac{3}{2^9}$, \emph{J. Algebra} \textbf{572} (2021),
60\textendash 75.

\bibitem[Zhu]{Z} Y. Zhu, Modular invariance of characters of vertex
operator algebras. \emph{J. Am. Math. Soc.} \textbf{9} (1996), 237\textendash 302.



\end{thebibliography}

\end{document}
